% Ngân hàng đề Toán 12 (KNTT) -- Chương II, Bài 8, Dạng 07
% Dạng: Tọa độ tổng, hiệu, tích với một số; độ dài vectơ, khoảng cách
%% Phương án ĐẦU TIÊN trong \cauTN luôn là đáp án đúng (gói trondethi tự đảo).

\cauTN[id=C2-B8-D07-NB-001,muc=NB,loigiai={Ta có $\overrightarrow{AB}=(x_B-x_A;y_B-y_A;z_B-z_A)=(2;0;0)$.}]{Trong không gian $Oxyz$, cho $A(-3;1;2)$, $B(-1;1;2)$. Tọa độ của vectơ $\overrightarrow{AB}$ là}
  {$(2;0;0)$.}
  {$(3;0;0)$.}{$(2;0;-1)$.}{$(-2;0;0)$.}

\cauTN[id=C2-B8-D07-NB-002,muc=NB,loigiai={Ta có $\overrightarrow{AB}=(x_B-x_A;y_B-y_A;z_B-z_A)=(-3;-1;0)$.}]{Trong không gian $Oxyz$, cho $A(2;2;-1)$, $B(-1;1;-1)$. Tọa độ của vectơ $\overrightarrow{AB}$ là}
  {$(-3;-1;0)$.}
  {$(-3;-1;1)$.}{$(-3;0;0)$.}{$(-3;1;0)$.}

\cauTN[id=C2-B8-D07-NB-003,muc=NB,loigiai={Ta có $\overrightarrow{AB}=(x_B-x_A;y_B-y_A;z_B-z_A)=(2;1;-3)$.}]{Trong không gian $Oxyz$, cho $A(-1;-1;1)$, $B(1;0;-2)$. Tọa độ của vectơ $\overrightarrow{AB}$ là}
  {$(2;1;-3)$.}
  {$(-2;1;-3)$.}{$(0;-1;-1)$.}{$(2;-3;1)$.}

\cauTN[id=C2-B8-D07-NB-004,muc=NB,loigiai={Ta có $\overrightarrow{AB}=(x_B-x_A;y_B-y_A;z_B-z_A)=(-3;-4;4)$.}]{Trong không gian $Oxyz$, cho $A(2;1;-1)$, $B(-1;-3;3)$. Tọa độ của vectơ $\overrightarrow{AB}$ là}
  {$(-3;-4;4)$.}
  {$(-3;-4;6)$.}{$(-3;4;4)$.}{$(-3;4;-4)$.}

\cauTN[id=C2-B8-D07-NB-005,muc=NB,loigiai={Ta có $\overrightarrow{AB}=(x_B-x_A;y_B-y_A;z_B-z_A)=(2;-4;6)$.}]{Trong không gian $Oxyz$, cho $A(-1;3;-3)$, $B(1;-1;3)$. Tọa độ của vectơ $\overrightarrow{AB}$ là}
  {$(2;-4;6)$.}
  {$(2;-3;6)$.}{$(2;-4;7)$.}{$(4;-4;6)$.}

\cauTN[id=C2-B8-D07-NB-006,muc=NB,loigiai={Ta có $\overrightarrow{AB}=(x_B-x_A;y_B-y_A;z_B-z_A)=(0;-4;-3)$.}]{Trong không gian $Oxyz$, cho $A(3;1;3)$, $B(3;-3;0)$. Tọa độ của vectơ $\overrightarrow{AB}$ là}
  {$(0;-4;-3)$.}
  {$(0;-4;-2)$.}{$(0;-6;-3)$.}{$(0;-3;-3)$.}

\cauTN[id=C2-B8-D07-NB-007,muc=NB,loigiai={Tính từng thành phần: $\vec{u}=-2\vec{a}+2\vec{b}=(-2;2;4)$.}]{Trong không gian $Oxyz$, cho $\vec{a}=(1;0;0)$, $\vec{b}=(0;1;2)$. Tọa độ của vectơ $\vec{u}=-2\vec{a}+2\vec{b}$ là}
  {$(-2;2;4)$.}
  {$(-2;-2;4)$.}{$(-2;2;2)$.}{$(-2;0;4)$.}

\cauTN[id=C2-B8-D07-NB-008,muc=NB,loigiai={Tính từng thành phần: $\vec{u}=3\vec{a}+\vec{b}-\vec{c}=(-11;0;0)$.}]{Trong không gian $Oxyz$, cho $\vec{a}=(-2;1;0)$, $\vec{b}=(-2;-2;-2)$, $\vec{c}=(3;1;-2)$. Tọa độ của vectơ $\vec{u}=3\vec{a}+\vec{b}-\vec{c}$ là}
  {$(-11;0;0)$.}
  {$(-11;0;-1)$.}{$(-10;0;0)$.}{$(-11;0;2)$.}

\cauTN[id=C2-B8-D07-NB-009,muc=NB,loigiai={Tính từng thành phần: $\vec{u}=3\vec{a}+2\vec{b}-\vec{c}=(13;0;9)$.}]{Trong không gian $Oxyz$, cho $\vec{a}=(3;1;3)$, $\vec{b}=(2;-3;0)$, $\vec{c}=(0;-3;0)$. Tọa độ của vectơ $\vec{u}=3\vec{a}+2\vec{b}-\vec{c}$ là}
  {$(13;0;9)$.}
  {$(12;0;9)$.}{$(14;0;9)$.}{$(13;1;9)$.}

\cauTN[id=C2-B8-D07-NB-010,muc=NB,loigiai={Tính từng thành phần: $\vec{u}=-2\vec{a}-\vec{b}-\vec{c}=(4;0;1)$.}]{Trong không gian $Oxyz$, cho $\vec{a}=(-3;0;-1)$, $\vec{b}=(0;1;-2)$, $\vec{c}=(2;-1;3)$. Tọa độ của vectơ $\vec{u}=-2\vec{a}-\vec{b}-\vec{c}$ là}
  {$(4;0;1)$.}
  {$(5;0;1)$.}{$(4;-2;1)$.}{$(6;0;1)$.}

\cauTN[id=C2-B8-D07-NB-011,muc=NB,loigiai={Tính từng thành phần: $\vec{u}=-\vec{a}-3\vec{b}+2\vec{c}=(12;8;-8)$.}]{Trong không gian $Oxyz$, cho $\vec{a}=(-1;3;-3)$, $\vec{b}=(-3;-3;3)$, $\vec{c}=(1;1;-1)$. Tọa độ của vectơ $\vec{u}=-\vec{a}-3\vec{b}+2\vec{c}$ là}
  {$(12;8;-8)$.}
  {$(12;8;-9)$.}{$(11;8;-8)$.}{$(14;8;-8)$.}

\cauTN[id=C2-B8-D07-NB-012,muc=NB,loigiai={Tính từng thành phần: $\vec{u}=2\vec{a}+\vec{b}=(6;4;-3)$.}]{Trong không gian $Oxyz$, cho $\vec{a}=(2;3;0)$, $\vec{b}=(2;-2;-3)$. Tọa độ của vectơ $\vec{u}=2\vec{a}+\vec{b}$ là}
  {$(6;4;-3)$.}
  {$(5;4;-3)$.}{$(6;5;-3)$.}{$(6;-3;4)$.}

\cauTN[id=C2-B8-D07-NB-013,muc=NB,loigiai={Tính từng thành phần: $\vec{u}=-2\vec{a}-3\vec{b}=(-3;3;-4)$.}]{Trong không gian $Oxyz$, cho $\vec{a}=(-3;-3;2)$, $\vec{b}=(3;1;0)$. Tọa độ của vectơ $\vec{u}=-2\vec{a}-3\vec{b}$ là}
  {$(-3;3;-4)$.}
  {$(-3;-4;3)$.}{$(3;3;-4)$.}{$(-3;5;-4)$.}

\cauTN[id=C2-B8-D07-NB-014,muc=NB,loigiai={Tính từng thành phần: $\vec{u}=-\vec{a}+\vec{b}=(0;-1;-1)$.}]{Trong không gian $Oxyz$, cho $\vec{a}=(1;-1;2)$, $\vec{b}=(1;-2;1)$. Tọa độ của vectơ $\vec{u}=-\vec{a}+\vec{b}$ là}
  {$(0;-1;-1)$.}
  {$(0;1;-1)$.}{$(0;0;-1)$.}{$(0;-3;-1)$.}

\cauTN[id=C2-B8-D07-NB-015,muc=NB,loigiai={Tính từng thành phần: $\vec{u}=3\vec{a}+3\vec{b}-\vec{c}=(0;-14;-3)$.}]{Trong không gian $Oxyz$, cho $\vec{a}=(-3;-2;2)$, $\vec{b}=(2;-3;-3)$, $\vec{c}=(-3;-1;0)$. Tọa độ của vectơ $\vec{u}=3\vec{a}+3\vec{b}-\vec{c}$ là}
  {$(0;-14;-3)$.}
  {$(0;-14;-4)$.}{$(0;-12;-3)$.}{$(0;14;3)$.}

\cauTN[id=C2-B8-D07-NB-016,muc=NB,loigiai={Tính từng thành phần: $\vec{u}=3\vec{a}-3\vec{b}=(6;0;-6)$.}]{Trong không gian $Oxyz$, cho $\vec{a}=(2;3;1)$, $\vec{b}=(0;3;3)$. Tọa độ của vectơ $\vec{u}=3\vec{a}-3\vec{b}$ là}
  {$(6;0;-6)$.}
  {$(6;-2;-6)$.}{$(6;0;-7)$.}{$(0;6;-6)$.}

\cauTN[id=C2-B8-D07-NB-017,muc=NB,loigiai={Ta có $|\vec{a}|=\sqrt{(-3)^2+(3)^2+(-2)^2}=\sqrt{22}=\sqrt{22}$.}]{Trong không gian $Oxyz$, cho $\vec{a}=(-3;3;-2)$. Độ dài của vectơ $\vec{a}$ bằng}
  {$\sqrt{22}$.}
  {$\sqrt{23}$.}{$\sqrt{21}$.}{$\sqrt{26}$.}

\cauTN[id=C2-B8-D07-NB-018,muc=NB,loigiai={Ta có $|\vec{a}|=\sqrt{(2)^2+(-2)^2+(2)^2}=\sqrt{12}=2\sqrt{3}$.}]{Trong không gian $Oxyz$, cho $\vec{a}=(2;-2;2)$. Độ dài của vectơ $\vec{a}$ bằng}
  {$2\sqrt{3}$.}
  {$\sqrt{13}$.}{$\sqrt{11}$.}{$4$.}

\cauTN[id=C2-B8-D07-NB-019,muc=NB,loigiai={Ta có $|\vec{a}|=\sqrt{(0)^2+(1)^2+(-1)^2}=\sqrt{2}=\sqrt{2}$.}]{Trong không gian $Oxyz$, cho $\vec{a}=(0;1;-1)$. Độ dài của vectơ $\vec{a}$ bằng}
  {$\sqrt{2}$.}
  {$\sqrt{3}$.}{$1$.}{$\sqrt{6}$.}

\cauTN[id=C2-B8-D07-NB-020,muc=NB,loigiai={Ta có $|\vec{a}|=\sqrt{(-2)^2+(-2)^2+(-2)^2}=\sqrt{12}=2\sqrt{3}$.}]{Trong không gian $Oxyz$, cho $\vec{a}=(-2;-2;-2)$. Độ dài của vectơ $\vec{a}$ bằng}
  {$2\sqrt{3}$.}
  {$\sqrt{13}$.}{$\sqrt{11}$.}{$4$.}

\cauTN[id=C2-B8-D07-NB-021,muc=NB,loigiai={Ta có $|\vec{a}|=\sqrt{(-1)^2+(-2)^2+(2)^2}=\sqrt{9}=3$.}]{Trong không gian $Oxyz$, cho $\vec{a}=(-1;-2;2)$. Độ dài của vectơ $\vec{a}$ bằng}
  {$3$.}
  {$\sqrt{10}$.}{$2\sqrt{2}$.}{$\sqrt{13}$.}

\cauTN[id=C2-B8-D07-NB-022,muc=NB,loigiai={Ta có $|\vec{a}|=\sqrt{(2)^2+(1)^2+(-2)^2}=\sqrt{9}=3$.}]{Trong không gian $Oxyz$, cho $\vec{a}=(2;1;-2)$. Độ dài của vectơ $\vec{a}$ bằng}
  {$3$.}
  {$\sqrt{10}$.}{$2\sqrt{2}$.}{$\sqrt{13}$.}

\cauTN[id=C2-B8-D07-NB-023,muc=NB,loigiai={Ta có $|\vec{a}|=\sqrt{(-1)^2+(2)^2+(1)^2}=\sqrt{6}=\sqrt{6}$.}]{Trong không gian $Oxyz$, cho $\vec{a}=(-1;2;1)$. Độ dài của vectơ $\vec{a}$ bằng}
  {$\sqrt{6}$.}
  {$\sqrt{7}$.}{$\sqrt{5}$.}{$\sqrt{10}$.}

\cauTN[id=C2-B8-D07-NB-024,muc=NB,loigiai={Ta có $|\vec{a}|=\sqrt{(3)^2+(-3)^2+(-2)^2}=\sqrt{22}=\sqrt{22}$.}]{Trong không gian $Oxyz$, cho $\vec{a}=(3;-3;-2)$. Độ dài của vectơ $\vec{a}$ bằng}
  {$\sqrt{22}$.}
  {$\sqrt{23}$.}{$\sqrt{21}$.}{$\sqrt{26}$.}

\cauTN[id=C2-B8-D07-NB-025,muc=NB,loigiai={Ta có $|\vec{a}|=\sqrt{(1)^2+(2)^2+(-3)^2}=\sqrt{14}=\sqrt{14}$.}]{Trong không gian $Oxyz$, cho $\vec{a}=(1;2;-3)$. Độ dài của vectơ $\vec{a}$ bằng}
  {$\sqrt{14}$.}
  {$\sqrt{15}$.}{$\sqrt{13}$.}{$3\sqrt{2}$.}

\cauTN[id=C2-B8-D07-NB-026,muc=NB,loigiai={Ta có $|\vec{a}|=\sqrt{(0)^2+(-1)^2+(-1)^2}=\sqrt{2}=\sqrt{2}$.}]{Trong không gian $Oxyz$, cho $\vec{a}=(0;-1;-1)$. Độ dài của vectơ $\vec{a}$ bằng}
  {$\sqrt{2}$.}
  {$\sqrt{3}$.}{$1$.}{$\sqrt{6}$.}

\cauTN[id=C2-B8-D07-NB-027,muc=NB,loigiai={Ta có $AB=\sqrt{(-4)^2+(2)^2+(4)^2}=\sqrt{36}=6$.}]{Trong không gian $Oxyz$, cho $A(1;0;-1)$, $B(-3;2;3)$. Khoảng cách giữa hai điểm $A$ và $B$ bằng}
  {$6$.}
  {$\sqrt{37}$.}{$\sqrt{35}$.}{$2\sqrt{10}$.}

\cauTN[id=C2-B8-D07-NB-028,muc=NB,loigiai={Ta có $AB=\sqrt{(3)^2+(4)^2+(2)^2}=\sqrt{29}=\sqrt{29}$.}]{Trong không gian $Oxyz$, cho $A(-3;-3;-3)$, $B(0;1;-1)$. Khoảng cách giữa hai điểm $A$ và $B$ bằng}
  {$\sqrt{29}$.}
  {$\sqrt{30}$.}{$2\sqrt{7}$.}{$\sqrt{33}$.}

\cauTN[id=C2-B8-D07-NB-029,muc=NB,loigiai={Ta có $AB=\sqrt{(-2)^2+(3)^2+(-5)^2}=\sqrt{38}=\sqrt{38}$.}]{Trong không gian $Oxyz$, cho $A(-1;-3;2)$, $B(-3;0;-3)$. Khoảng cách giữa hai điểm $A$ và $B$ bằng}
  {$\sqrt{38}$.}
  {$\sqrt{39}$.}{$\sqrt{37}$.}{$\sqrt{42}$.}

\cauTN[id=C2-B8-D07-NB-030,muc=NB,loigiai={Ta có $AB=\sqrt{(-1)^2+(2)^2+(-2)^2}=\sqrt{9}=3$.}]{Trong không gian $Oxyz$, cho $A(-1;0;3)$, $B(-2;2;1)$. Khoảng cách giữa hai điểm $A$ và $B$ bằng}
  {$3$.}
  {$\sqrt{10}$.}{$2\sqrt{2}$.}{$\sqrt{13}$.}

\cauTN[id=C2-B8-D07-NB-031,muc=NB,loigiai={Ta có $AB=\sqrt{(-3)^2+(-5)^2+(-2)^2}=\sqrt{38}=\sqrt{38}$.}]{Trong không gian $Oxyz$, cho $A(2;2;2)$, $B(-1;-3;0)$. Khoảng cách giữa hai điểm $A$ và $B$ bằng}
  {$\sqrt{38}$.}
  {$\sqrt{39}$.}{$\sqrt{37}$.}{$\sqrt{42}$.}

\cauTN[id=C2-B8-D07-NB-032,muc=NB,loigiai={Ta có $AB=\sqrt{(-5)^2+(3)^2+(-1)^2}=\sqrt{35}=\sqrt{35}$.}]{Trong không gian $Oxyz$, cho $A(3;-2;-1)$, $B(-2;1;-2)$. Khoảng cách giữa hai điểm $A$ và $B$ bằng}
  {$\sqrt{35}$.}
  {$6$.}{$\sqrt{34}$.}{$\sqrt{39}$.}

\cauTN[id=C2-B8-D07-NB-033,muc=NB,loigiai={Ta có $AB=\sqrt{(2)^2+(-3)^2+(-1)^2}=\sqrt{14}=\sqrt{14}$.}]{Trong không gian $Oxyz$, cho $A(0;3;1)$, $B(2;0;0)$. Khoảng cách giữa hai điểm $A$ và $B$ bằng}
  {$\sqrt{14}$.}
  {$\sqrt{15}$.}{$\sqrt{13}$.}{$3\sqrt{2}$.}

\cauTN[id=C2-B8-D07-NB-034,muc=NB,loigiai={Ta có $AB=\sqrt{(-2)^2+(2)^2+(3)^2}=\sqrt{17}=\sqrt{17}$.}]{Trong không gian $Oxyz$, cho $A(0;-3;-3)$, $B(-2;-1;0)$. Khoảng cách giữa hai điểm $A$ và $B$ bằng}
  {$\sqrt{17}$.}
  {$3\sqrt{2}$.}{$4$.}{$\sqrt{21}$.}

\cauTN[id=C2-B8-D07-NB-035,muc=NB,loigiai={Ta có $AB=\sqrt{(2)^2+(-1)^2+(-2)^2}=\sqrt{9}=3$.}]{Trong không gian $Oxyz$, cho $A(-1;0;0)$, $B(1;-1;-2)$. Khoảng cách giữa hai điểm $A$ và $B$ bằng}
  {$3$.}
  {$\sqrt{10}$.}{$2\sqrt{2}$.}{$\sqrt{13}$.}

\cauTN[id=C2-B8-D07-NB-036,muc=NB,loigiai={Ta có $AB=\sqrt{(1)^2+(4)^2+(4)^2}=\sqrt{33}=\sqrt{33}$.}]{Trong không gian $Oxyz$, cho $A(-3;-3;-3)$, $B(-2;1;1)$. Khoảng cách giữa hai điểm $A$ và $B$ bằng}
  {$\sqrt{33}$.}
  {$\sqrt{34}$.}{$4\sqrt{2}$.}{$\sqrt{37}$.}

\cauDS[id=C2-B8-D07-NB-037,muc=NB,loigiai={Kiểm tra từng mệnh đề bằng công thức tọa độ. a) Sai. b) Sai. c) Đúng. d) Đúng.}]{Trong không gian $Oxyz$, cho $A(2;2;3)$, $B(2;1;3)$.  Cho $\vec{a}=(-1;-2;-2)$, $\vec{b}=(0;-3;-1)$, $\vec{c}=(0;1;3)$. Xét tính đúng sai của các mệnh đề sau:}{%
  \y{Độ dài đoạn thẳng $AB$ bằng $\sqrt{5}$.}{s}
  \y{Độ dài của vectơ $\vec{a}$ bằng $\sqrt{10}$.}{s}
  \y{Vectơ $\vec{u}=3\vec{a}+2\vec{b}$ có tọa độ $(-3;-12;-8)$.}{d}
  \y{Vectơ $\vec{u}=-2\vec{a}-\vec{b}$ có tọa độ $(2;7;5)$.}{d}}

\cauDS[id=C2-B8-D07-NB-038,muc=NB,loigiai={Kiểm tra từng mệnh đề bằng công thức tọa độ. a) Sai. b) Sai. c) Sai. d) Đúng.}]{Trong không gian $Oxyz$, cho $A(-2;2;2)$, $B(0;0;-2)$.  Cho $\vec{a}=(-2;0;0)$, $\vec{b}=(-1;-3;0)$, $\vec{c}=(-2;1;-1)$. Xét tính đúng sai của các mệnh đề sau:}{%
  \y{Độ dài của vectơ $\vec{a}$ bằng $\sqrt{3}$.}{s}
  \y{Vectơ $\vec{u}=-2\vec{a}+3\vec{b}+\vec{c}$ có tọa độ $(-1;-8;1)$.}{s}
  \y{Vectơ $\vec{u}=3\vec{a}+\vec{b}+\vec{c}$ có tọa độ $(-9;-3;-1)$.}{s}
  \y{Độ dài đoạn thẳng $AB$ bằng $2\sqrt{6}$.}{d}}

\cauDS[id=C2-B8-D07-NB-039,muc=NB,loigiai={Kiểm tra từng mệnh đề bằng công thức tọa độ. a) Đúng. b) Sai. c) Đúng. d) Sai.}]{Trong không gian $Oxyz$, cho $A(3;3;-3)$, $B(3;3;2)$.  Cho $\vec{a}=(2;0;-3)$, $\vec{b}=(0;0;2)$, $\vec{c}=(-1;-3;-2)$. Xét tính đúng sai của các mệnh đề sau:}{%
  \y{Độ dài của vectơ $\vec{a}$ bằng $\sqrt{13}$.}{d}
  \y{Độ dài của vectơ $\vec{a}$ bằng $\sqrt{14}$.}{s}
  \y{Độ dài đoạn thẳng $AB$ bằng $5$.}{d}
  \y{Vectơ $\vec{u}=2\vec{a}+2\vec{b}+2\vec{c}$ có tọa độ $(2;-6;-4)$.}{s}}

\cauDS[id=C2-B8-D07-NB-040,muc=NB,loigiai={Kiểm tra từng mệnh đề bằng công thức tọa độ. a) Sai. b) Sai. c) Đúng. d) Đúng.}]{Trong không gian $Oxyz$, cho $A(-2;3;3)$, $B(3;-3;0)$.  Cho $\vec{a}=(2;3;-3)$, $\vec{b}=(-2;-3;1)$, $\vec{c}=(-1;-2;-1)$. Xét tính đúng sai của các mệnh đề sau:}{%
  \y{Độ dài đoạn thẳng $AB$ bằng $\sqrt{74}$.}{s}
  \y{Vectơ $\vec{u}=-\vec{a}-3\vec{b}+\vec{c}$ có tọa độ $(3;4;0)$.}{s}
  \y{Vectơ $\overrightarrow{AB}$ có tọa độ $(5;-6;-3)$.}{d}
  \y{Vectơ $\vec{u}=2\vec{a}+2\vec{b}$ có tọa độ $(0;0;-4)$.}{d}}

\cauDS[id=C2-B8-D07-NB-041,muc=NB,loigiai={Kiểm tra từng mệnh đề bằng công thức tọa độ. a) Sai. b) Đúng. c) Đúng. d) Sai.}]{Trong không gian $Oxyz$, cho $A(0;2;-2)$, $B(2;-3;1)$.  Cho $\vec{a}=(1;-1;3)$, $\vec{b}=(0;2;-3)$, $\vec{c}=(3;-2;-3)$. Xét tính đúng sai của các mệnh đề sau:}{%
  \y{Vectơ $\overrightarrow{AB}$ có tọa độ $(0;-5;3)$.}{s}
  \y{Vectơ $\vec{u}=2\vec{a}-3\vec{b}$ có tọa độ $(2;-8;15)$.}{d}
  \y{Độ dài của vectơ $\vec{a}$ bằng $\sqrt{11}$.}{d}
  \y{Độ dài của vectơ $\vec{a}$ bằng $\sqrt{15}$.}{s}}

\cauDS[id=C2-B8-D07-NB-042,muc=NB,loigiai={Kiểm tra từng mệnh đề bằng công thức tọa độ. a) Đúng. b) Sai. c) Đúng. d) Sai.}]{Trong không gian $Oxyz$, cho $A(-3;0;2)$, $B(3;0;3)$.  Cho $\vec{a}=(2;-1;1)$, $\vec{b}=(-3;2;0)$, $\vec{c}=(3;-1;1)$. Xét tính đúng sai của các mệnh đề sau:}{%
  \y{Độ dài của vectơ $\vec{a}$ bằng $\sqrt{6}$.}{d}
  \y{Vectơ $\vec{u}=2\vec{a}+2\vec{b}-\vec{c}$ có tọa độ $(5;-3;-1)$.}{s}
  \y{Vectơ $\overrightarrow{AB}$ có tọa độ $(6;0;1)$.}{d}
  \y{Độ dài của vectơ $\vec{a}$ bằng $\sqrt{5}$.}{s}}

\cauDS[id=C2-B8-D07-NB-043,muc=NB,loigiai={Kiểm tra từng mệnh đề bằng công thức tọa độ. a) Đúng. b) Sai. c) Sai. d) Đúng.}]{Trong không gian $Oxyz$, cho $A(1;1;2)$, $B(-3;1;1)$.  Cho $\vec{a}=(-1;-2;-3)$, $\vec{b}=(1;0;-2)$, $\vec{c}=(2;2;-1)$. Xét tính đúng sai của các mệnh đề sau:}{%
  \y{Vectơ $\vec{u}=2\vec{a}-3\vec{b}+2\vec{c}$ có tọa độ $(-1;0;-2)$.}{d}
  \y{Vectơ $\vec{u}=3\vec{a}-\vec{b}+\vec{c}$ có tọa độ $(-2;-4;-7)$.}{s}
  \y{Độ dài của vectơ $\vec{a}$ bằng $\sqrt{15}$.}{s}
  \y{Độ dài đoạn thẳng $AB$ bằng $\sqrt{17}$.}{d}}

\cauDS[id=C2-B8-D07-NB-044,muc=NB,loigiai={Kiểm tra từng mệnh đề bằng công thức tọa độ. a) Đúng. b) Đúng. c) Sai. d) Sai.}]{Trong không gian $Oxyz$, cho $A(2;-1;2)$, $B(-2;-1;3)$.  Cho $\vec{a}=(-3;-1;-1)$, $\vec{b}=(-3;2;2)$, $\vec{c}=(-2;-3;3)$. Xét tính đúng sai của các mệnh đề sau:}{%
  \y{Độ dài đoạn thẳng $AB$ bằng $\sqrt{17}$.}{d}
  \y{Vectơ $\vec{u}=2\vec{a}+3\vec{b}$ có tọa độ $(-15;4;4)$.}{d}
  \y{Độ dài của vectơ $\vec{a}$ bằng $\sqrt{10}$.}{s}
  \y{Vectơ $\overrightarrow{AB}$ có tọa độ $(-4;0;0)$.}{s}}

\cauDS[id=C2-B8-D07-NB-045,muc=NB,loigiai={Kiểm tra từng mệnh đề bằng công thức tọa độ. a) Sai. b) Sai. c) Đúng. d) Đúng.}]{Trong không gian $Oxyz$, cho $A(1;0;-3)$, $B(-2;0;3)$.  Cho $\vec{a}=(-1;0;-2)$, $\vec{b}=(3;2;-2)$, $\vec{c}=(-3;-2;2)$. Xét tính đúng sai của các mệnh đề sau:}{%
  \y{Vectơ $\vec{u}=3\vec{a}+\vec{b}$ có tọa độ $(0;1;-8)$.}{s}
  \y{Độ dài đoạn thẳng $AB$ bằng $\sqrt{46}$.}{s}
  \y{Độ dài của vectơ $\vec{a}$ bằng $\sqrt{5}$.}{d}
  \y{Độ dài đoạn thẳng $AB$ bằng $3\sqrt{5}$.}{d}}

\cauDS[id=C2-B8-D07-NB-046,muc=NB,loigiai={Kiểm tra từng mệnh đề bằng công thức tọa độ. a) Đúng. b) Sai. c) Sai. d) Sai.}]{Trong không gian $Oxyz$, cho $A(3;3;-1)$, $B(-2;2;-3)$.  Cho $\vec{a}=(-1;-1;-2)$. Xét tính đúng sai của các mệnh đề sau:}{%
  \y{Vectơ $\overrightarrow{AB}$ có tọa độ $(-5;-1;-2)$.}{d}
  \y{Độ dài đoạn thẳng $AB$ bằng $\sqrt{34}$.}{s}
  \y{Độ dài của vectơ $\vec{a}$ bằng $\sqrt{5}$.}{s}
  \y{Vectơ $\overrightarrow{AB}$ có tọa độ $(-5;-1;-1)$.}{s}}

\cauDS[id=C2-B8-D07-NB-047,muc=NB,loigiai={Kiểm tra từng mệnh đề bằng công thức tọa độ. a) Sai. b) Đúng. c) Sai. d) Đúng.}]{Trong không gian $Oxyz$, cho $A(2;-3;1)$, $B(-1;-2;-3)$.  Cho $\vec{a}=(-1;1;-2)$, $\vec{b}=(2;-3;2)$, $\vec{c}=(3;-1;0)$. Xét tính đúng sai của các mệnh đề sau:}{%
  \y{Vectơ $\overrightarrow{AB}$ có tọa độ $(-4;1;-4)$.}{s}
  \y{Vectơ $\overrightarrow{AB}$ có tọa độ $(-3;1;-4)$.}{d}
  \y{Vectơ $\vec{u}=-2\vec{a}-\vec{b}+2\vec{c}$ có tọa độ $(6;-1;0)$.}{s}
  \y{Độ dài đoạn thẳng $AB$ bằng $\sqrt{26}$.}{d}}

\cauDS[id=C2-B8-D07-NB-048,muc=NB,loigiai={Kiểm tra từng mệnh đề bằng công thức tọa độ. a) Sai. b) Sai. c) Sai. d) Đúng.}]{Trong không gian $Oxyz$, cho $A(2;2;1)$, $B(0;1;-1)$.  Cho $\vec{a}=(3;-3;-1)$. Xét tính đúng sai của các mệnh đề sau:}{%
  \y{Vectơ $\overrightarrow{AB}$ có tọa độ $(2;-1;-2)$.}{s}
  \y{Độ dài đoạn thẳng $AB$ bằng $\sqrt{10}$.}{s}
  \y{Độ dài đoạn thẳng $AB$ bằng $2\sqrt{2}$.}{s}
  \y{Độ dài của vectơ $\vec{a}$ bằng $\sqrt{19}$.}{d}}

\cauTLN[id=C2-B8-D07-NB-049,muc=NB,loigiai={Tính từng thành phần: $\vec{u}=-\vec{a}-\vec{b}=(0;-1;-3)$.}]{Trong không gian $Oxyz$, cho $\vec{a}=(1;-1;0)$, $\vec{b}=(-1;2;3)$. Tính tổng các tọa độ của vectơ $\vec{u}=-\vec{a}-\vec{b}$.}{-4}

\cauTLN[id=C2-B8-D07-NB-050,muc=NB,loigiai={Tính từng thành phần: $\vec{u}=-\vec{a}-3\vec{b}-\vec{c}=(-6;4;-2)$.}]{Trong không gian $Oxyz$, cho $\vec{a}=(0;-1;2)$, $\vec{b}=(3;-1;-1)$, $\vec{c}=(-3;0;3)$. Tính tổng các tọa độ của vectơ $\vec{u}=-\vec{a}-3\vec{b}-\vec{c}$.}{-4}

\cauTLN[id=C2-B8-D07-NB-051,muc=NB,loigiai={Tính từng thành phần: $\vec{u}=3\vec{a}+2\vec{b}=(12;-6;3)$.}]{Trong không gian $Oxyz$, cho $\vec{a}=(2;0;3)$, $\vec{b}=(3;-3;-3)$. Tính tổng các tọa độ của vectơ $\vec{u}=3\vec{a}+2\vec{b}$.}{9}

\cauTLN[id=C2-B8-D07-NB-052,muc=NB,loigiai={Tính từng thành phần: $\vec{u}=3\vec{a}+3\vec{b}=(0;0;-3)$.}]{Trong không gian $Oxyz$, cho $\vec{a}=(0;1;-3)$, $\vec{b}=(0;-1;2)$. Tính tổng các tọa độ của vectơ $\vec{u}=3\vec{a}+3\vec{b}$.}{-3}

\cauTLN[id=C2-B8-D07-NB-053,muc=NB,loigiai={Ta có $|\vec{a}|=\sqrt{(3)^2+(-1)^2+(2)^2}=\sqrt{14}=\sqrt{14}$.}]{Trong không gian $Oxyz$, cho $\vec{a}=(3;-1;2)$. Tính $|\vec{a}|^2$.}{14}

\cauTLN[id=C2-B8-D07-NB-054,muc=NB,loigiai={Ta có $|\vec{a}|=\sqrt{(3)^2+(-2)^2+(-3)^2}=\sqrt{22}=\sqrt{22}$.}]{Trong không gian $Oxyz$, cho $\vec{a}=(3;-2;-3)$. Tính $|\vec{a}|^2$.}{22}

\cauTLN[id=C2-B8-D07-NB-055,muc=NB,loigiai={Ta có $|\vec{a}|=\sqrt{(-1)^2+(-2)^2+(-2)^2}=\sqrt{9}=3$.}]{Trong không gian $Oxyz$, cho $\vec{a}=(-1;-2;-2)$. Tính $|\vec{a}|^2$.}{9}

\cauTLN[id=C2-B8-D07-NB-056,muc=NB,loigiai={Ta có $|\vec{a}|=\sqrt{(1)^2+(-2)^2+(2)^2}=\sqrt{9}=3$.}]{Trong không gian $Oxyz$, cho $\vec{a}=(1;-2;2)$. Tính $|\vec{a}|^2$.}{9}

\cauTLN[id=C2-B8-D07-NB-057,muc=NB,loigiai={Ta có $AB=\sqrt{(4)^2+(0)^2+(-3)^2}=\sqrt{25}=5$.}]{Trong không gian $Oxyz$, cho $A(-1;-1;2)$, $B(3;-1;-1)$. Tính $AB^2$.}{25}

\cauTLN[id=C2-B8-D07-NB-058,muc=NB,loigiai={Ta có $AB=\sqrt{(0)^2+(2)^2+(6)^2}=\sqrt{40}=2\sqrt{10}$.}]{Trong không gian $Oxyz$, cho $A(1;1;-3)$, $B(1;3;3)$. Tính $AB^2$.}{40}

\cauTLN[id=C2-B8-D07-NB-059,muc=NB,loigiai={Ta có $AB=\sqrt{(3)^2+(4)^2+(3)^2}=\sqrt{34}=\sqrt{34}$.}]{Trong không gian $Oxyz$, cho $A(-2;-1;-3)$, $B(1;3;0)$. Tính $AB^2$.}{34}

\cauTLN[id=C2-B8-D07-NB-060,muc=NB,loigiai={Ta có $AB=\sqrt{(2)^2+(-1)^2+(0)^2}=\sqrt{5}=\sqrt{5}$.}]{Trong không gian $Oxyz$, cho $A(-2;2;-1)$, $B(0;1;-1)$. Tính $AB^2$.}{5}

\cauTN[id=C2-B8-D07-TH-001,muc=TH,loigiai={Ta có $|\vec{a}|=3$ nên vectơ đơn vị cùng hướng với $\vec{a}$ là $\dfrac{1}{|\vec{a}|}\vec{a}=(-\dfrac{2}{3};-\dfrac{1}{3};-\dfrac{2}{3})$.}]{Vectơ đơn vị $\vec{e}$ cùng hướng với vectơ $\vec{a}=(-2;-1;-2)$ có tọa độ là}
  {$(-\dfrac{2}{3};-\dfrac{1}{3};-\dfrac{2}{3})$.}
  {$(-\dfrac{2}{3};-\dfrac{1}{3};-\dfrac{8}{3})$.}{$(-\dfrac{2}{3};-\dfrac{1}{3};\dfrac{4}{3})$.}{$(-\dfrac{2}{3};-\dfrac{2}{3};-\dfrac{1}{3})$.}

\cauTN[id=C2-B8-D07-TH-002,muc=TH,loigiai={Ta có $|\vec{a}|=9$ nên vectơ đơn vị cùng hướng với $\vec{a}$ là $\dfrac{1}{|\vec{a}|}\vec{a}=(\dfrac{4}{9};\dfrac{4}{9};\dfrac{7}{9})$.}]{Vectơ đơn vị $\vec{e}$ cùng hướng với vectơ $\vec{a}=(4;4;7)$ có tọa độ là}
  {$(\dfrac{4}{9};\dfrac{4}{9};\dfrac{7}{9})$.}
  {$(-\dfrac{4}{9};-\dfrac{4}{9};-\dfrac{7}{9})$.}{$(-\dfrac{14}{9};\dfrac{4}{9};\dfrac{7}{9})$.}{$(\dfrac{4}{9};-\dfrac{4}{9};\dfrac{7}{9})$.}

\cauTN[id=C2-B8-D07-TH-003,muc=TH,loigiai={Ta có $|\vec{a}|=5$ nên vectơ đơn vị cùng hướng với $\vec{a}$ là $\dfrac{1}{|\vec{a}|}\vec{a}=(0;\dfrac{3}{5};\dfrac{4}{5})$.}]{Vectơ đơn vị $\vec{e}$ cùng hướng với vectơ $\vec{a}=(0;3;4)$ có tọa độ là}
  {$(0;\dfrac{3}{5};\dfrac{4}{5})$.}
  {$(-1;\dfrac{3}{5};\dfrac{4}{5})$.}{$(2;\dfrac{3}{5};\dfrac{4}{5})$.}{$(0;-\dfrac{3}{5};\dfrac{4}{5})$.}

\cauTN[id=C2-B8-D07-TH-004,muc=TH,loigiai={Ta có $|\vec{a}|=3$ nên vectơ đơn vị cùng hướng với $\vec{a}$ là $\dfrac{1}{|\vec{a}|}\vec{a}=(-\dfrac{2}{3};\dfrac{2}{3};\dfrac{1}{3})$.}]{Vectơ đơn vị $\vec{e}$ cùng hướng với vectơ $\vec{a}=(-2;2;1)$ có tọa độ là}
  {$(-\dfrac{2}{3};\dfrac{2}{3};\dfrac{1}{3})$.}
  {$(\dfrac{4}{3};\dfrac{2}{3};\dfrac{1}{3})$.}{$(\dfrac{2}{3};-\dfrac{2}{3};-\dfrac{1}{3})$.}{$(-\dfrac{2}{3};-\dfrac{1}{3};\dfrac{1}{3})$.}

\cauTN[id=C2-B8-D07-TH-005,muc=TH,loigiai={Ta có $|\vec{a}|=10$ nên vectơ đơn vị cùng hướng với $\vec{a}$ là $\dfrac{1}{|\vec{a}|}\vec{a}=(-\dfrac{3}{5};0;-\dfrac{4}{5})$.}]{Vectơ đơn vị $\vec{e}$ cùng hướng với vectơ $\vec{a}=(-6;0;-8)$ có tọa độ là}
  {$(-\dfrac{3}{5};0;-\dfrac{4}{5})$.}
  {$(\dfrac{3}{5};0;\dfrac{4}{5})$.}{$(-\dfrac{8}{5};0;-\dfrac{4}{5})$.}{$(-\dfrac{3}{5};-\dfrac{4}{5};0)$.}

\cauTN[id=C2-B8-D07-TH-006,muc=TH,loigiai={Ta có $|\vec{a}|=2$ nên vectơ đơn vị cùng hướng với $\vec{a}$ là $\dfrac{1}{|\vec{a}|}\vec{a}=(0;0;1)$.}]{Vectơ đơn vị $\vec{e}$ cùng hướng với vectơ $\vec{a}=(0;0;2)$ có tọa độ là}
  {$(0;0;1)$.}
  {$(-1;0;1)$.}{$(0;2;1)$.}{$(0;0;\dfrac{1}{2})$.}

\cauTN[id=C2-B8-D07-TH-007,muc=TH,loigiai={Ta có $|\vec{a}|=9$ nên vectơ đơn vị cùng hướng với $\vec{a}$ là $\dfrac{1}{|\vec{a}|}\vec{a}=(-\dfrac{4}{9};-\dfrac{4}{9};-\dfrac{7}{9})$.}]{Vectơ đơn vị $\vec{e}$ cùng hướng với vectơ $\vec{a}=(-4;-4;-7)$ có tọa độ là}
  {$(-\dfrac{4}{9};-\dfrac{4}{9};-\dfrac{7}{9})$.}
  {$(-\dfrac{4}{9};-\dfrac{4}{9};\dfrac{11}{9})$.}{$(-\dfrac{4}{9};-\dfrac{4}{9};\dfrac{7}{9})$.}{$(-\dfrac{7}{9};-\dfrac{4}{9};-\dfrac{4}{9})$.}

\cauTN[id=C2-B8-D07-TH-008,muc=TH,loigiai={Ta có $|\vec{a}|=10$ nên vectơ đơn vị cùng hướng với $\vec{a}$ là $\dfrac{1}{|\vec{a}|}\vec{a}=(\dfrac{4}{5};0;\dfrac{3}{5})$.}]{Vectơ đơn vị $\vec{e}$ cùng hướng với vectơ $\vec{a}=(8;0;6)$ có tọa độ là}
  {$(\dfrac{4}{5};0;\dfrac{3}{5})$.}
  {$(\dfrac{4}{5};0;-\dfrac{3}{5})$.}{$(\dfrac{4}{5};-2;\dfrac{3}{5})$.}{$(\dfrac{4}{5};1;\dfrac{3}{5})$.}

\cauTN[id=C2-B8-D07-TH-009,muc=TH,loigiai={Ta có $|\vec{a}|=10$ nên vectơ đơn vị cùng hướng với $\vec{a}$ là $\dfrac{1}{|\vec{a}|}\vec{a}=(-1;0;0)$.}]{Vectơ đơn vị $\vec{e}$ cùng hướng với vectơ $\vec{a}=(-10;0;0)$ có tọa độ là}
  {$(-1;0;0)$.}
  {$(-1;0;-1)$.}{$(-1;2;0)$.}{$(0;-1;0)$.}

\cauTN[id=C2-B8-D07-TH-010,muc=TH,loigiai={Ta có $|\vec{a}|=1$ nên vectơ đơn vị cùng hướng với $\vec{a}$ là $\dfrac{1}{|\vec{a}|}\vec{a}=(0;0;1)$.}]{Vectơ đơn vị $\vec{e}$ cùng hướng với vectơ $\vec{a}=(0;0;1)$ có tọa độ là}
  {$(0;0;1)$.}
  {$(0;0;-1)$.}{$(1;0;1)$.}{$(2;0;1)$.}

\cauTN[id=C2-B8-D07-TH-011,muc=TH,loigiai={Ta có $|\vec{a}|=11$ nên vectơ đơn vị cùng hướng với $\vec{a}$ là $\dfrac{1}{|\vec{a}|}\vec{a}=(-\dfrac{6}{11};-\dfrac{6}{11};-\dfrac{7}{11})$.}]{Vectơ đơn vị $\vec{e}$ cùng hướng với vectơ $\vec{a}=(-6;-6;-7)$ có tọa độ là}
  {$(-\dfrac{6}{11};-\dfrac{6}{11};-\dfrac{7}{11})$.}
  {$(\dfrac{5}{11};-\dfrac{6}{11};-\dfrac{7}{11})$.}{$(-\dfrac{6}{121};-\dfrac{6}{121};-\dfrac{7}{121})$.}{$(\dfrac{6}{11};\dfrac{6}{11};\dfrac{7}{11})$.}

\cauTN[id=C2-B8-D07-TH-012,muc=TH,loigiai={Ta có $|\vec{a}|=5$ nên vectơ đơn vị cùng hướng với $\vec{a}$ là $\dfrac{1}{|\vec{a}|}\vec{a}=(0;\dfrac{4}{5};\dfrac{3}{5})$.}]{Vectơ đơn vị $\vec{e}$ cùng hướng với vectơ $\vec{a}=(0;4;3)$ có tọa độ là}
  {$(0;\dfrac{4}{5};\dfrac{3}{5})$.}
  {$(\dfrac{3}{5};\dfrac{4}{5};0)$.}{$(0;\dfrac{4}{5};-\dfrac{2}{5})$.}{$(\dfrac{4}{5};0;\dfrac{3}{5})$.}

\cauTN[id=C2-B8-D07-TH-013,muc=TH,loigiai={Ta có $B=A+\overrightarrow{AB}=(-4;3;5)+((1;-1;-1))=(-3;2;4)$.}]{Cho điểm $A(-4;3;5)$, $\vec{a}=(-1;0;2)$, $\vec{b}=(-3;1;5)$. Điểm $B$ thoả mãn $\overrightarrow{AB}=2\vec{a}-\vec{b}$ có tọa độ là}
  {$(-3;2;4)$.}
  {$(-3;2;6)$.}{$(-3;3;4)$.}{$(3;-2;-4)$.}

\cauTN[id=C2-B8-D07-TH-014,muc=TH,loigiai={Ta có $B=A+\overrightarrow{AB}=(5;-4;-2)+((13;-7;1))=(18;-11;-1)$.}]{Cho điểm $A(5;-4;-2)$, $\vec{a}=(5;-1;-1)$, $\vec{b}=(1;2;-2)$. Điểm $B$ thoả mãn $\overrightarrow{AB}=3\vec{a}-2\vec{b}$ có tọa độ là}
  {$(18;-11;-1)$.}
  {$(18;-11;1)$.}{$(20;-11;-1)$.}{$(18;-11;-3)$.}

\cauTN[id=C2-B8-D07-TH-015,muc=TH,loigiai={Ta có $B=A+\overrightarrow{AB}=(-4;-4;-2)+((4;-8;0))=(0;-12;-2)$.}]{Cho điểm $A(-4;-4;-2)$, $\vec{a}=(0;-4;-5)$, $\vec{b}=(-2;0;-5)$. Điểm $B$ thoả mãn $\overrightarrow{AB}=2\vec{a}-2\vec{b}$ có tọa độ là}
  {$(0;-12;-2)$.}
  {$(-2;-12;0)$.}{$(2;-12;-2)$.}{$(0;12;-2)$.}

\cauTN[id=C2-B8-D07-TH-016,muc=TH,loigiai={Ta có $B=A+\overrightarrow{AB}=(-1;-4;2)+((-3;12;4))=(-4;8;6)$.}]{Cho điểm $A(-1;-4;2)$, $\vec{a}=(-3;4;0)$, $\vec{b}=(0;4;2)$. Điểm $B$ thoả mãn $\overrightarrow{AB}=\vec{a}+2\vec{b}$ có tọa độ là}
  {$(-4;8;6)$.}
  {$(-4;8;4)$.}{$(-4;-8;6)$.}{$(-6;8;6)$.}

\cauTN[id=C2-B8-D07-TH-017,muc=TH,loigiai={Ta có $B=A+\overrightarrow{AB}=(2;3;-3)+((-1;-6;6))=(1;-3;3)$.}]{Cho điểm $A(2;3;-3)$, $\vec{a}=(2;2;-4)$, $\vec{b}=(1;-4;2)$. Điểm $B$ thoả mãn $\overrightarrow{AB}=-1\vec{a}+\vec{b}$ có tọa độ là}
  {$(1;-3;3)$.}
  {$(1;-3;1)$.}{$(0;-3;3)$.}{$(1;-5;3)$.}

\cauTN[id=C2-B8-D07-TH-018,muc=TH,loigiai={Ta có $B=A+\overrightarrow{AB}=(0;-4;5)+((-8;-10;-5))=(-8;-14;0)$.}]{Cho điểm $A(0;-4;5)$, $\vec{a}=(-4;-4;1)$, $\vec{b}=(2;3;3)$. Điểm $B$ thoả mãn $\overrightarrow{AB}=\vec{a}-2\vec{b}$ có tọa độ là}
  {$(-8;-14;0)$.}
  {$(8;-14;0)$.}{$(-8;-15;0)$.}{$(-8;-12;0)$.}

\cauTN[id=C2-B8-D07-TH-019,muc=TH,loigiai={Ta có $B=A+\overrightarrow{AB}=(-2;0;-5)+((20;4;-14))=(18;4;-19)$.}]{Cho điểm $A(-2;0;-5)$, $\vec{a}=(5;2;-5)$, $\vec{b}=(5;0;-2)$. Điểm $B$ thoả mãn $\overrightarrow{AB}=2\vec{a}+2\vec{b}$ có tọa độ là}
  {$(18;4;-19)$.}
  {$(18;4;-17)$.}{$(16;4;-19)$.}{$(-18;-4;19)$.}

\cauTN[id=C2-B8-D07-TH-020,muc=TH,loigiai={Ta có $B=A+\overrightarrow{AB}=(2;4;0)+((5;-7;-7))=(7;-3;-7)$.}]{Cho điểm $A(2;4;0)$, $\vec{a}=(3;-3;1)$, $\vec{b}=(-1;2;4)$. Điểm $B$ thoả mãn $\overrightarrow{AB}=\vec{a}-2\vec{b}$ có tọa độ là}
  {$(7;-3;-7)$.}
  {$(7;-5;-7)$.}{$(7;-3;-8)$.}{$(8;-3;-7)$.}

\cauTN[id=C2-B8-D07-TH-021,muc=TH,loigiai={Ta có $B=A+\overrightarrow{AB}=(3;-3;-4)+((8;6;10))=(11;3;6)$.}]{Cho điểm $A(3;-3;-4)$, $\vec{a}=(2;-4;0)$, $\vec{b}=(3;5;5)$. Điểm $B$ thoả mãn $\overrightarrow{AB}=\vec{a}+2\vec{b}$ có tọa độ là}
  {$(11;3;6)$.}
  {$(11;3;7)$.}{$(11;4;6)$.}{$(11;3;8)$.}

\cauTN[id=C2-B8-D07-TH-022,muc=TH,loigiai={Ta có $B=A+\overrightarrow{AB}=(-4;-3;-1)+((-13;-4;-1))=(-17;-7;-2)$.}]{Cho điểm $A(-4;-3;-1)$, $\vec{a}=(-5;2;3)$, $\vec{b}=(-4;-3;-2)$. Điểm $B$ thoả mãn $\overrightarrow{AB}=\vec{a}+2\vec{b}$ có tọa độ là}
  {$(-17;-7;-2)$.}
  {$(17;7;2)$.}{$(-16;-7;-2)$.}{$(-15;-7;-2)$.}

\cauTN[id=C2-B8-D07-TH-023,muc=TH,loigiai={Ta có $B=A+\overrightarrow{AB}=(-4;-2;2)+((-1;-2;-7))=(-5;-4;-5)$.}]{Cho điểm $A(-4;-2;2)$, $\vec{a}=(3;1;-4)$, $\vec{b}=(4;3;3)$. Điểm $B$ thoả mãn $\overrightarrow{AB}=\vec{a}-\vec{b}$ có tọa độ là}
  {$(-5;-4;-5)$.}
  {$(-5;-5;-4)$.}{$(-6;-4;-5)$.}{$(5;-4;-5)$.}

\cauTN[id=C2-B8-D07-TH-024,muc=TH,loigiai={Ta có $B=A+\overrightarrow{AB}=(-3;-4;-5)+((-2;0;-7))=(-5;-4;-12)$.}]{Cho điểm $A(-3;-4;-5)$, $\vec{a}=(-2;-2;-1)$, $\vec{b}=(2;3;-2)$. Điểm $B$ thoả mãn $\overrightarrow{AB}=3\vec{a}+2\vec{b}$ có tọa độ là}
  {$(-5;-4;-12)$.}
  {$(-5;-2;-12)$.}{$(-1;-4;2)$.}{$(-5;-4;-11)$.}

\cauTN[id=C2-B8-D07-TH-025,muc=TH,loigiai={Ta có $-2\vec{a}-3\vec{b}=(5;21;7)$ nên $|-2\vec{a}-3\vec{b}|=\sqrt{515}=\sqrt{515}$.}]{Cho $\vec{a}=(-1;-3;-5)$, $\vec{b}=(-1;-5;1)$. Độ dài của vectơ $-2\vec{a}-3\vec{b}$ bằng}
  {$\sqrt{515}$.}
  {$2\sqrt{129}$.}{$\sqrt{514}$.}{$\sqrt{519}$.}

\cauTN[id=C2-B8-D07-TH-026,muc=TH,loigiai={Ta có $2\vec{a}+\vec{b}=(-4;-4;2)$ nên $|2\vec{a}+\vec{b}|=\sqrt{36}=6$.}]{Cho $\vec{a}=(0;-2;1)$, $\vec{b}=(-4;0;0)$. Độ dài của vectơ $2\vec{a}+\vec{b}$ bằng}
  {$6$.}
  {$\sqrt{37}$.}{$\sqrt{35}$.}{$2\sqrt{10}$.}

\cauTN[id=C2-B8-D07-TH-027,muc=TH,loigiai={Ta có $-\vec{a}+3\vec{b}=(2;17;1)$ nên $|-\vec{a}+3\vec{b}|=\sqrt{294}=7\sqrt{6}$.}]{Cho $\vec{a}=(1;-2;5)$, $\vec{b}=(1;5;2)$. Độ dài của vectơ $-\vec{a}+3\vec{b}$ bằng}
  {$7\sqrt{6}$.}
  {$\sqrt{295}$.}{$\sqrt{293}$.}{$\sqrt{298}$.}

\cauTN[id=C2-B8-D07-TH-028,muc=TH,loigiai={Ta có $-2\vec{a}-\vec{b}=(8;-4;0)$ nên $|-2\vec{a}-\vec{b}|=\sqrt{80}=4\sqrt{5}$.}]{Cho $\vec{a}=(-2;0;-2)$, $\vec{b}=(-4;4;4)$. Độ dài của vectơ $-2\vec{a}-\vec{b}$ bằng}
  {$4\sqrt{5}$.}
  {$9$.}{$\sqrt{79}$.}{$2\sqrt{21}$.}

\cauTN[id=C2-B8-D07-TH-029,muc=TH,loigiai={Ta có $3\vec{a}+3\vec{b}=(0;9;0)$ nên $|3\vec{a}+3\vec{b}|=\sqrt{81}=9$.}]{Cho $\vec{a}=(-2;2;5)$, $\vec{b}=(2;1;-5)$. Độ dài của vectơ $3\vec{a}+3\vec{b}$ bằng}
  {$9$.}
  {$\sqrt{82}$.}{$4\sqrt{5}$.}{$\sqrt{85}$.}

\cauTN[id=C2-B8-D07-TH-030,muc=TH,loigiai={Ta có $-\vec{a}+3\vec{b}=(7;4;15)$ nên $|-\vec{a}+3\vec{b}|=\sqrt{290}=\sqrt{290}$.}]{Cho $\vec{a}=(-1;-1;0)$, $\vec{b}=(2;1;5)$. Độ dài của vectơ $-\vec{a}+3\vec{b}$ bằng}
  {$\sqrt{290}$.}
  {$\sqrt{291}$.}{$17$.}{$7\sqrt{6}$.}

\cauTN[id=C2-B8-D07-TH-031,muc=TH,loigiai={Ta có $-\vec{a}-\vec{b}=(-5;-2;4)$ nên $|-\vec{a}-\vec{b}|=\sqrt{45}=3\sqrt{5}$.}]{Cho $\vec{a}=(2;1;1)$, $\vec{b}=(3;1;-5)$. Độ dài của vectơ $-\vec{a}-\vec{b}$ bằng}
  {$3\sqrt{5}$.}
  {$\sqrt{46}$.}{$2\sqrt{11}$.}{$7$.}

\cauTN[id=C2-B8-D07-TH-032,muc=TH,loigiai={Ta có $-2\vec{a}+2\vec{b}=(-2;2;-2)$ nên $|-2\vec{a}+2\vec{b}|=\sqrt{12}=2\sqrt{3}$.}]{Cho $\vec{a}=(1;-5;-3)$, $\vec{b}=(0;-4;-4)$. Độ dài của vectơ $-2\vec{a}+2\vec{b}$ bằng}
  {$2\sqrt{3}$.}
  {$\sqrt{13}$.}{$\sqrt{11}$.}{$4$.}

\cauTN[id=C2-B8-D07-TH-033,muc=TH,loigiai={Ta có $-2\vec{a}+3\vec{b}=(2;19;-2)$ nên $|-2\vec{a}+3\vec{b}|=\sqrt{369}=3\sqrt{41}$.}]{Cho $\vec{a}=(5;-5;-2)$, $\vec{b}=(4;3;-2)$. Độ dài của vectơ $-2\vec{a}+3\vec{b}$ bằng}
  {$3\sqrt{41}$.}
  {$\sqrt{370}$.}{$4\sqrt{23}$.}{$\sqrt{373}$.}

\cauTN[id=C2-B8-D07-TH-034,muc=TH,loigiai={Ta có $3\vec{a}-3\vec{b}=(3;-6;0)$ nên $|3\vec{a}-3\vec{b}|=\sqrt{45}=3\sqrt{5}$.}]{Cho $\vec{a}=(5;2;-2)$, $\vec{b}=(4;4;-2)$. Độ dài của vectơ $3\vec{a}-3\vec{b}$ bằng}
  {$3\sqrt{5}$.}
  {$\sqrt{46}$.}{$2\sqrt{11}$.}{$7$.}

\cauTN[id=C2-B8-D07-TH-035,muc=TH,loigiai={Ta có $-2\vec{a}-\vec{b}=(8;1;1)$ nên $|-2\vec{a}-\vec{b}|=\sqrt{66}=\sqrt{66}$.}]{Cho $\vec{a}=(-3;-2;-1)$, $\vec{b}=(-2;3;1)$. Độ dài của vectơ $-2\vec{a}-\vec{b}$ bằng}
  {$\sqrt{66}$.}
  {$\sqrt{67}$.}{$\sqrt{65}$.}{$\sqrt{70}$.}

\cauTN[id=C2-B8-D07-TH-036,muc=TH,loigiai={Ta có $3\vec{a}-3\vec{b}=(27;-9;-3)$ nên $|3\vec{a}-3\vec{b}|=\sqrt{819}=3\sqrt{91}$.}]{Cho $\vec{a}=(4;0;-2)$, $\vec{b}=(-5;3;-1)$. Độ dài của vectơ $3\vec{a}-3\vec{b}$ bằng}
  {$3\sqrt{91}$.}
  {$2\sqrt{205}$.}{$\sqrt{818}$.}{$\sqrt{823}$.}

\cauTN[id=C2-B8-D07-TH-037,muc=TH,loigiai={Tính từng thành phần: $\vec{u}=2\vec{a}-3\vec{b}-\vec{c}=(23;-17;-1)$.}]{Trong không gian $Oxyz$, cho $\vec{a}=(3;-4;5)$, $\vec{b}=(-5;3;5)$, $\vec{c}=(-2;0;-4)$. Tọa độ của vectơ $\vec{u}=2\vec{a}-3\vec{b}-\vec{c}$ là}
  {$(23;-17;-1)$.}
  {$(23;-18;-1)$.}{$(23;-17;-3)$.}{$(24;-17;-1)$.}

\cauTN[id=C2-B8-D07-TH-038,muc=TH,loigiai={Tính từng thành phần: $\vec{u}=3\vec{a}-\vec{b}+2\vec{c}=(4;-13;20)$.}]{Trong không gian $Oxyz$, cho $\vec{a}=(0;-4;3)$, $\vec{b}=(2;-3;-1)$, $\vec{c}=(3;-2;5)$. Tọa độ của vectơ $\vec{u}=3\vec{a}-\vec{b}+2\vec{c}$ là}
  {$(4;-13;20)$.}
  {$(-4;-13;20)$.}{$(4;13;20)$.}{$(5;-13;20)$.}

\cauTN[id=C2-B8-D07-TH-039,muc=TH,loigiai={Tính từng thành phần: $\vec{u}=3\vec{a}+2\vec{b}+2\vec{c}=(12;7;8)$.}]{Trong không gian $Oxyz$, cho $\vec{a}=(2;1;2)$, $\vec{b}=(3;0;1)$, $\vec{c}=(0;2;0)$. Tọa độ của vectơ $\vec{u}=3\vec{a}+2\vec{b}+2\vec{c}$ là}
  {$(12;7;8)$.}
  {$(12;7;7)$.}{$(11;7;8)$.}{$(12;6;8)$.}

\cauTN[id=C2-B8-D07-TH-040,muc=TH,loigiai={Tính từng thành phần: $\vec{u}=3\vec{a}+2\vec{b}=(-14;-9;5)$.}]{Trong không gian $Oxyz$, cho $\vec{a}=(-4;-3;-1)$, $\vec{b}=(-1;0;4)$. Tọa độ của vectơ $\vec{u}=3\vec{a}+2\vec{b}$ là}
  {$(-14;-9;5)$.}
  {$(-13;-9;5)$.}{$(-14;-9;6)$.}{$(-9;-14;5)$.}

\cauTN[id=C2-B8-D07-TH-041,muc=TH,loigiai={Tính từng thành phần: $\vec{u}=-2\vec{a}-\vec{b}+2\vec{c}=(13;5;7)$.}]{Trong không gian $Oxyz$, cho $\vec{a}=(-5;1;1)$, $\vec{b}=(-1;-1;-5)$, $\vec{c}=(1;3;2)$. Tọa độ của vectơ $\vec{u}=-2\vec{a}-\vec{b}+2\vec{c}$ là}
  {$(13;5;7)$.}
  {$(13;-5;7)$.}{$(11;5;7)$.}{$(13;5;8)$.}

\cauTN[id=C2-B8-D07-TH-042,muc=TH,loigiai={Tính từng thành phần: $\vec{u}=-2\vec{a}-\vec{b}=(2;-10;10)$.}]{Trong không gian $Oxyz$, cho $\vec{a}=(-1;4;-5)$, $\vec{b}=(0;2;0)$. Tọa độ của vectơ $\vec{u}=-2\vec{a}-\vec{b}$ là}
  {$(2;-10;10)$.}
  {$(2;-9;10)$.}{$(2;-10;12)$.}{$(2;-10;-10)$.}

\cauTN[id=C2-B8-D07-TH-043,muc=TH,loigiai={Tính từng thành phần: $\vec{u}=3\vec{a}-3\vec{b}-\vec{c}=(2;30;2)$.}]{Trong không gian $Oxyz$, cho $\vec{a}=(-3;5;0)$, $\vec{b}=(-4;-5;-2)$, $\vec{c}=(1;0;4)$. Tọa độ của vectơ $\vec{u}=3\vec{a}-3\vec{b}-\vec{c}$ là}
  {$(2;30;2)$.}
  {$(2;30;4)$.}{$(-2;-30;-2)$.}{$(2;29;2)$.}

\cauTN[id=C2-B8-D07-TH-044,muc=TH,loigiai={Tính từng thành phần: $\vec{u}=-\vec{a}-3\vec{b}=(17;10;10)$.}]{Trong không gian $Oxyz$, cho $\vec{a}=(-2;5;-1)$, $\vec{b}=(-5;-5;-3)$. Tọa độ của vectơ $\vec{u}=-\vec{a}-3\vec{b}$ là}
  {$(17;10;10)$.}
  {$(17;11;10)$.}{$(17;10;8)$.}{$(17;10;9)$.}

\cauTN[id=C2-B8-D07-TH-045,muc=TH,loigiai={Tính từng thành phần: $\vec{u}=-2\vec{a}+2\vec{b}=(-14;-2;-2)$.}]{Trong không gian $Oxyz$, cho $\vec{a}=(2;-2;-3)$, $\vec{b}=(-5;-3;-4)$. Tọa độ của vectơ $\vec{u}=-2\vec{a}+2\vec{b}$ là}
  {$(-14;-2;-2)$.}
  {$(-15;-2;-2)$.}{$(-14;-2;-3)$.}{$(-14;-2;2)$.}

\cauTN[id=C2-B8-D07-TH-046,muc=TH,loigiai={Tính từng thành phần: $\vec{u}=-2\vec{a}+2\vec{b}+2\vec{c}=(4;8;-10)$.}]{Trong không gian $Oxyz$, cho $\vec{a}=(-2;-3;5)$, $\vec{b}=(3;-2;-1)$, $\vec{c}=(-3;3;1)$. Tọa độ của vectơ $\vec{u}=-2\vec{a}+2\vec{b}+2\vec{c}$ là}
  {$(4;8;-10)$.}
  {$(5;8;-10)$.}{$(4;-8;-10)$.}{$(2;8;-10)$.}

\cauTN[id=C2-B8-D07-TH-047,muc=TH,loigiai={Tính từng thành phần: $\vec{u}=-2\vec{a}-\vec{b}=(10;-2;1)$.}]{Trong không gian $Oxyz$, cho $\vec{a}=(-4;0;0)$, $\vec{b}=(-2;2;-1)$. Tọa độ của vectơ $\vec{u}=-2\vec{a}-\vec{b}$ là}
  {$(10;-2;1)$.}
  {$(9;-2;1)$.}{$(-10;-2;1)$.}{$(10;-1;1)$.}

\cauTN[id=C2-B8-D07-TH-048,muc=TH,loigiai={Tính từng thành phần: $\vec{u}=-\vec{a}+\vec{b}-\vec{c}=(11;-2;2)$.}]{Trong không gian $Oxyz$, cho $\vec{a}=(-5;1;-2)$, $\vec{b}=(5;4;0)$, $\vec{c}=(-1;5;0)$. Tọa độ của vectơ $\vec{u}=-\vec{a}+\vec{b}-\vec{c}$ là}
  {$(11;-2;2)$.}
  {$(9;-2;2)$.}{$(-11;-2;2)$.}{$(12;-2;2)$.}

\cauDS[id=C2-B8-D07-TH-049,muc=TH,loigiai={Kiểm tra từng mệnh đề bằng công thức tọa độ. a) Sai. b) Đúng. c) Đúng. d) Đúng.}]{Trong không gian $Oxyz$, cho $A(-5;-2;-2)$, $B(3;5;2)$.  Cho $\vec{a}=(-5;-3;0)$, $\vec{b}=(3;-5;-5)$, $\vec{c}=(4;-2;2)$. Xét tính đúng sai của các mệnh đề sau:}{%
  \y{Độ dài đoạn thẳng $AB$ bằng $\sqrt{133}$.}{s}
  \y{Với $\vec{a}=(-5;-3;0)$, $\vec{b}=(3;-5;-5)$ thì $|2\vec{a}-\vec{b}|=\sqrt{195}$.}{d}
  \y{Vectơ $\vec{u}=-\vec{a}-3\vec{b}-\vec{c}$ có tọa độ $(-8;20;13)$.}{d}
  \y{Độ dài của vectơ $\vec{a}$ bằng $\sqrt{34}$.}{d}}

\cauDS[id=C2-B8-D07-TH-050,muc=TH,loigiai={Kiểm tra từng mệnh đề bằng công thức tọa độ. a) Đúng. b) Sai. c) Sai. d) Sai.}]{Trong không gian $Oxyz$, cho $A(-3;1;-1)$, $B(-1;2;-1)$.  Cho $\vec{a}=(-5;2;3)$. Xét tính đúng sai của các mệnh đề sau:}{%
  \y{Với $\vec{a}=(-5;2;3)$, $\vec{b}=(-4;-1;-1)$ thì $|-\vec{a}+3\vec{b}|=\sqrt{110}$.}{d}
  \y{Độ dài đoạn thẳng $AB$ bằng $\sqrt{6}$.}{s}
  \y{Với $\vec{a}=(-5;2;3)$, $\vec{b}=(-4;-1;-1)$ thì $|2\vec{a}+2\vec{b}|=7\sqrt{7}$.}{s}
  \y{Độ dài của vectơ $\vec{a}$ bằng $\sqrt{42}$.}{s}}

\cauDS[id=C2-B8-D07-TH-051,muc=TH,loigiai={Kiểm tra từng mệnh đề bằng công thức tọa độ. a) Đúng. b) Sai. c) Đúng. d) Sai.}]{Trong không gian $Oxyz$, cho $\vec{a}=(0;-4;1)$, $\vec{b}=(5;1;0)$, $\vec{c}=(4;-3;1)$. Xét tính đúng sai của các mệnh đề sau:}{%
  \y{Độ dài của vectơ $\vec{a}$ bằng $\sqrt{17}$.}{d}
  \y{Độ dài của vectơ $\vec{a}$ bằng $4$.}{s}
  \y{Vectơ $\vec{u}=-\vec{a}-3\vec{b}$ có tọa độ $(-15;1;-1)$.}{d}
  \y{Với $\vec{a}=(0;-4;1)$, $\vec{b}=(5;1;0)$ thì $|3\vec{a}-\vec{b}|=3\sqrt{23}$.}{s}}

\cauDS[id=C2-B8-D07-TH-052,muc=TH,loigiai={Kiểm tra từng mệnh đề bằng công thức tọa độ. a) Sai. b) Sai. c) Đúng. d) Sai.}]{Trong không gian $Oxyz$, cho $A(-5;0;5)$, $B(-2;4;-3)$.  Cho $\vec{a}=(-3;-4;3)$, $\vec{b}=(-2;-3;0)$, $\vec{c}=(-1;1;2)$. Xét tính đúng sai của các mệnh đề sau:}{%
  \y{Độ dài đoạn thẳng $AB$ bằng $\sqrt{93}$.}{s}
  \y{Vectơ $\vec{u}=-\vec{a}+2\vec{b}$ có tọa độ $(-3;-2;-1)$.}{s}
  \y{Với $\vec{a}=(-3;-4;3)$, $\vec{b}=(-2;-3;0)$ thì $|-2\vec{a}-3\vec{b}|=\sqrt{469}$.}{d}
  \y{Độ dài của vectơ $\vec{a}$ bằng $\sqrt{33}$.}{s}}

\cauDS[id=C2-B8-D07-TH-053,muc=TH,loigiai={Kiểm tra từng mệnh đề bằng công thức tọa độ. a) Sai. b) Đúng. c) Sai. d) Sai.}]{Trong không gian $Oxyz$, cho $A(4;5;-3)$, $B(-4;0;-5)$.  Cho $\vec{a}=(-1;-3;-1)$, $\vec{b}=(-3;3;-1)$, $\vec{c}=(1;-5;0)$. Xét tính đúng sai của các mệnh đề sau:}{%
  \y{Độ dài đoạn thẳng $AB$ bằng $\sqrt{97}$.}{s}
  \y{Độ dài đoạn thẳng $AB$ bằng $\sqrt{93}$.}{d}
  \y{Với $\vec{a}=(-1;-3;-1)$, $\vec{b}=(-3;3;-1)$ thì $|3\vec{a}-3\vec{b}|=\sqrt{359}$.}{s}
  \y{Vectơ $\vec{u}=2\vec{a}-3\vec{b}-\vec{c}$ có tọa độ $(6;-10;-1)$.}{s}}

\cauDS[id=C2-B8-D07-TH-054,muc=TH,loigiai={Kiểm tra từng mệnh đề bằng công thức tọa độ. a) Đúng. b) Đúng. c) Sai. d) Sai.}]{Trong không gian $Oxyz$, cho $A(3;5;1)$, $B(-2;5;0)$.  Cho $\vec{a}=(1;0;-5)$, $\vec{b}=(-5;3;5)$, $\vec{c}=(-2;-4;0)$. Xét tính đúng sai của các mệnh đề sau:}{%
  \y{Độ dài của vectơ $\vec{a}$ bằng $\sqrt{26}$.}{d}
  \y{Với $\vec{a}=(1;0;-5)$, $\vec{b}=(-5;3;5)$ thì $|2\vec{a}-3\vec{b}|=\sqrt{995}$.}{d}
  \y{Độ dài đoạn thẳng $AB$ bằng $\sqrt{30}$.}{s}
  \y{Vectơ $\vec{u}=-2\vec{a}-3\vec{b}$ có tọa độ $(13;-5;-9)$.}{s}}

\cauDS[id=C2-B8-D07-TH-055,muc=TH,loigiai={Kiểm tra từng mệnh đề bằng công thức tọa độ. a) Đúng. b) Đúng. c) Đúng. d) Sai.}]{Trong không gian $Oxyz$, cho $A(4;3;4)$, $B(-1;1;-2)$.  Cho $\vec{a}=(1;2;-5)$, $\vec{b}=(0;5;-5)$, $\vec{c}=(0;1;4)$. Xét tính đúng sai của các mệnh đề sau:}{%
  \y{Độ dài của vectơ $\vec{a}$ bằng $\sqrt{30}$.}{d}
  \y{Độ dài đoạn thẳng $AB$ bằng $\sqrt{65}$.}{d}
  \y{Vectơ $\vec{u}=3\vec{a}+\vec{b}-\vec{c}$ có tọa độ $(3;10;-24)$.}{d}
  \y{Với $\vec{a}=(1;2;-5)$, $\vec{b}=(0;5;-5)$ thì $|2\vec{a}+\vec{b}|=\sqrt{311}$.}{s}}

\cauDS[id=C2-B8-D07-TH-056,muc=TH,loigiai={Kiểm tra từng mệnh đề bằng công thức tọa độ. a) Đúng. b) Sai. c) Đúng. d) Sai.}]{Trong không gian $Oxyz$, cho $A(3;-2;-5)$, $B(-3;-3;1)$.  Cho $\vec{a}=(5;3;5)$, $\vec{b}=(3;3;4)$, $\vec{c}=(4;4;-3)$. Xét tính đúng sai của các mệnh đề sau:}{%
  \y{Với $\vec{a}=(5;3;5)$, $\vec{b}=(3;3;4)$ thì $|3\vec{a}-\vec{b}|=\sqrt{301}$.}{d}
  \y{Độ dài đoạn thẳng $AB$ bằng $6\sqrt{2}$.}{s}
  \y{Độ dài đoạn thẳng $AB$ bằng $\sqrt{73}$.}{d}
  \y{Vectơ $\vec{u}=-2\vec{a}+3\vec{b}+\vec{c}$ có tọa độ $(3;9;-1)$.}{s}}

\cauDS[id=C2-B8-D07-TH-057,muc=TH,loigiai={Kiểm tra từng mệnh đề bằng công thức tọa độ. a) Đúng. b) Đúng. c) Đúng. d) Sai.}]{Trong không gian $Oxyz$, cho $\vec{a}=(-5;-3;2)$, $\vec{b}=(5;-4;0)$, $\vec{c}=(0;3;-1)$. Xét tính đúng sai của các mệnh đề sau:}{%
  \y{Với $\vec{a}=(-5;-3;2)$, $\vec{b}=(5;-4;0)$ thì $|-\vec{a}+3\vec{b}|=\sqrt{485}$.}{d}
  \y{Độ dài của vectơ $\vec{a}$ bằng $\sqrt{38}$.}{d}
  \y{Vectơ $\vec{u}=-\vec{a}+\vec{b}$ có tọa độ $(10;-1;-2)$.}{d}
  \y{Với $\vec{a}=(-5;-3;2)$, $\vec{b}=(5;-4;0)$ thì $|3\vec{a}-3\vec{b}|=\sqrt{949}$.}{s}}

\cauDS[id=C2-B8-D07-TH-058,muc=TH,loigiai={Kiểm tra từng mệnh đề bằng công thức tọa độ. a) Đúng. b) Đúng. c) Sai. d) Sai.}]{Trong không gian $Oxyz$, cho $A(3;-3;-3)$, $B(5;0;-1)$.  Cho $\vec{a}=(-3;3;-1)$. Xét tính đúng sai của các mệnh đề sau:}{%
  \y{Độ dài đoạn thẳng $AB$ bằng $\sqrt{17}$.}{d}
  \y{Với $\vec{a}=(-3;3;-1)$, $\vec{b}=(1;-3;3)$ thì $|-2\vec{a}+2\vec{b}|=4\sqrt{17}$.}{d}
  \y{Độ dài của vectơ $\vec{a}$ bằng $3\sqrt{2}$.}{s}
  \y{Với $\vec{a}=(-3;3;-1)$, $\vec{b}=(1;-3;3)$ thì $|-\vec{a}+3\vec{b}|=2\sqrt{71}$.}{s}}

\cauDS[id=C2-B8-D07-TH-059,muc=TH,loigiai={Kiểm tra từng mệnh đề bằng công thức tọa độ. a) Đúng. b) Đúng. c) Sai. d) Đúng.}]{Trong không gian $Oxyz$, cho $A(-5;-1;-4)$, $B(-1;-1;3)$.  Cho $\vec{a}=(3;-3;3)$. Xét tính đúng sai của các mệnh đề sau:}{%
  \y{Độ dài của vectơ $\vec{a}$ bằng $3\sqrt{3}$.}{d}
  \y{Với $\vec{a}=(3;-3;3)$, $\vec{b}=(1;-4;-1)$ thì $|-2\vec{a}+\vec{b}|=\sqrt{78}$.}{d}
  \y{Với $\vec{a}=(3;-3;3)$, $\vec{b}=(1;-4;-1)$ thì $|3\vec{a}+2\vec{b}|=\sqrt{463}$.}{s}
  \y{Độ dài đoạn thẳng $AB$ bằng $\sqrt{65}$.}{d}}

\cauDS[id=C2-B8-D07-TH-060,muc=TH,loigiai={Kiểm tra từng mệnh đề bằng công thức tọa độ. a) Đúng. b) Sai. c) Đúng. d) Sai.}]{Trong không gian $Oxyz$, cho $A(1;0;4)$, $B(-1;1;0)$.  Cho $\vec{a}=(-4;0;3)$, $\vec{b}=(4;-1;3)$, $\vec{c}=(-1;2;-5)$. Xét tính đúng sai của các mệnh đề sau:}{%
  \y{Độ dài đoạn thẳng $AB$ bằng $\sqrt{21}$.}{d}
  \y{Với $\vec{a}=(-4;0;3)$, $\vec{b}=(4;-1;3)$ thì $|-\vec{a}-\vec{b}|=\sqrt{38}$.}{s}
  \y{Vectơ $\vec{u}=-\vec{a}-\vec{b}-\vec{c}$ có tọa độ $(1;-1;-1)$.}{d}
  \y{Độ dài của vectơ $\vec{a}$ bằng $\sqrt{29}$.}{s}}

\cauDS[id=C2-B8-D07-TH-061,muc=TH,loigiai={Kiểm tra từng mệnh đề bằng công thức tọa độ. a) Sai. b) Sai. c) Sai. d) Đúng.}]{Trong không gian $Oxyz$, cho $A(-4;1;-1)$, $B(-4;-3;-2)$.  Cho $\vec{a}=(4;3;2)$. Xét tính đúng sai của các mệnh đề sau:}{%
  \y{Với $\vec{a}=(4;3;2)$, $\vec{b}=(-3;-1;0)$ thì $|2\vec{a}-3\vec{b}|=\sqrt{385}$.}{s}
  \y{Độ dài đoạn thẳng $AB$ bằng $3\sqrt{2}$.}{s}
  \y{Độ dài của vectơ $\vec{a}$ bằng $\sqrt{33}$.}{s}
  \y{Với $\vec{a}=(4;3;2)$, $\vec{b}=(-3;-1;0)$ thì $|3\vec{a}-\vec{b}|=19$.}{d}}

\cauDS[id=C2-B8-D07-TH-062,muc=TH,loigiai={Kiểm tra từng mệnh đề bằng công thức tọa độ. a) Sai. b) Đúng. c) Sai. d) Sai.}]{Trong không gian $Oxyz$, cho $A(-5;2;-2)$, $B(-4;-5;0)$.  Cho $\vec{a}=(1;2;4)$. Xét tính đúng sai của các mệnh đề sau:}{%
  \y{Độ dài đoạn thẳng $AB$ bằng $\sqrt{53}$.}{s}
  \y{Độ dài của vectơ $\vec{a}$ bằng $\sqrt{21}$.}{d}
  \y{Độ dài của vectơ $\vec{a}$ bằng $5$.}{s}
  \y{Với $\vec{a}=(1;2;4)$, $\vec{b}=(-4;4;1)$ thì $|-\vec{a}+2\vec{b}|=5\sqrt{5}$.}{s}}

\cauDS[id=C2-B8-D07-TH-063,muc=TH,loigiai={Kiểm tra từng mệnh đề bằng công thức tọa độ. a) Đúng. b) Đúng. c) Đúng. d) Sai.}]{Trong không gian $Oxyz$, cho $A(5;-3;5)$, $B(-1;-1;0)$.  Cho $\vec{a}=(-3;-5;-5)$. Xét tính đúng sai của các mệnh đề sau:}{%
  \y{Với $\vec{a}=(-3;-5;-5)$, $\vec{b}=(-3;3;1)$ thì $|3\vec{a}-\vec{b}|=2\sqrt{154}$.}{d}
  \y{Độ dài đoạn thẳng $AB$ bằng $\sqrt{65}$.}{d}
  \y{Độ dài của vectơ $\vec{a}$ bằng $\sqrt{59}$.}{d}
  \y{Độ dài của vectơ $\vec{a}$ bằng $\sqrt{58}$.}{s}}

\cauDS[id=C2-B8-D07-TH-064,muc=TH,loigiai={Kiểm tra từng mệnh đề bằng công thức tọa độ. a) Đúng. b) Đúng. c) Sai. d) Đúng.}]{Trong không gian $Oxyz$, cho $A(5;-2;1)$, $B(-3;-3;5)$.  Cho $\vec{a}=(-4;-3;0)$. Xét tính đúng sai của các mệnh đề sau:}{%
  \y{Độ dài đoạn thẳng $AB$ bằng $9$.}{d}
  \y{Độ dài của vectơ $\vec{a}$ bằng $5$.}{d}
  \y{Độ dài của vectơ $\vec{a}$ bằng $\sqrt{29}$.}{s}
  \y{Với $\vec{a}=(-4;-3;0)$, $\vec{b}=(-2;2;1)$ thì $|-\vec{a}+3\vec{b}|=\sqrt{94}$.}{d}}

\cauTLN[id=C2-B8-D07-TH-065,muc=TH,loigiai={Ta có $B=A+\overrightarrow{AB}=(1;4;5)+((7;4;1))=(8;8;6)$.}]{Cho điểm $A(1;4;5)$, $\vec{a}=(3;0;5)$, $\vec{b}=(-5;-2;-3)$. Điểm $B$ thoả mãn $\overrightarrow{AB}=-1\vec{a}-2\vec{b}$. Tính $x_B+y_B+z_B$.}{22}

\cauTLN[id=C2-B8-D07-TH-066,muc=TH,loigiai={Ta có $B=A+\overrightarrow{AB}=(5;5;-3)+((13;-10;-2))=(18;-5;-5)$.}]{Cho điểm $A(5;5;-3)$, $\vec{a}=(-5;0;-4)$, $\vec{b}=(4;-5;-3)$. Điểm $B$ thoả mãn $\overrightarrow{AB}=-1\vec{a}+2\vec{b}$. Tính $x_B+y_B+z_B$.}{8}

\cauTLN[id=C2-B8-D07-TH-067,muc=TH,loigiai={Ta có $B=A+\overrightarrow{AB}=(5;-1;4)+((10;-2;0))=(15;-3;4)$.}]{Cho điểm $A(5;-1;4)$, $\vec{a}=(0;-1;3)$, $\vec{b}=(5;0;-3)$. Điểm $B$ thoả mãn $\overrightarrow{AB}=2\vec{a}+2\vec{b}$. Tính $x_B+y_B+z_B$.}{16}

\cauTLN[id=C2-B8-D07-TH-068,muc=TH,loigiai={Ta có $B=A+\overrightarrow{AB}=(-3;4;-1)+((-16;8;-2))=(-19;12;-3)$.}]{Cho điểm $A(-3;4;-1)$, $\vec{a}=(-4;-1;-5)$, $\vec{b}=(4;-5;-4)$. Điểm $B$ thoả mãn $\overrightarrow{AB}=2\vec{a}-2\vec{b}$. Tính $x_B+y_B+z_B$.}{-10}

\cauTLN[id=C2-B8-D07-TH-069,muc=TH,loigiai={Ta có $B=A+\overrightarrow{AB}=(-1;-2;0)+((19;-20;7))=(18;-22;7)$.}]{Cho điểm $A(-1;-2;0)$, $\vec{a}=(3;-4;-1)$, $\vec{b}=(-5;4;-5)$. Điểm $B$ thoả mãn $\overrightarrow{AB}=3\vec{a}-2\vec{b}$. Tính $x_B+y_B+z_B$.}{3}

\cauTLN[id=C2-B8-D07-TH-070,muc=TH,loigiai={Ta có $B=A+\overrightarrow{AB}=(-2;4;-4)+((-7;3;3))=(-9;7;-1)$.}]{Cho điểm $A(-2;4;-4)$, $\vec{a}=(-4;1;0)$, $\vec{b}=(1;1;3)$. Điểm $B$ thoả mãn $\overrightarrow{AB}=2\vec{a}+\vec{b}$. Tính $x_B+y_B+z_B$.}{-3}

\cauTLN[id=C2-B8-D07-TH-071,muc=TH,loigiai={Ta có $2\vec{a}+3\vec{b}=(11;-13;12)$ nên $|2\vec{a}+3\vec{b}|=\sqrt{434}=\sqrt{434}$.}]{Cho $\vec{a}=(1;-5;0)$, $\vec{b}=(3;-1;4)$. Tính $|2\vec{a}+3\vec{b}|^2$.}{434}

\cauTLN[id=C2-B8-D07-TH-072,muc=TH,loigiai={Ta có $-\vec{a}+\vec{b}=(0;-9;-1)$ nên $|-\vec{a}+\vec{b}|=\sqrt{82}=\sqrt{82}$.}]{Cho $\vec{a}=(-2;4;4)$, $\vec{b}=(-2;-5;3)$. Tính $|-\vec{a}+\vec{b}|^2$.}{82}

\cauTLN[id=C2-B8-D07-TH-073,muc=TH,loigiai={Ta có $-2\vec{a}+2\vec{b}=(-6;6;6)$ nên $|-2\vec{a}+2\vec{b}|=\sqrt{108}=6\sqrt{3}$.}]{Cho $\vec{a}=(4;-3;-3)$, $\vec{b}=(1;0;0)$. Tính $|-2\vec{a}+2\vec{b}|^2$.}{108}

\cauTLN[id=C2-B8-D07-TH-074,muc=TH,loigiai={Ta có $-\vec{a}+\vec{b}=(2;-6;1)$ nên $|-\vec{a}+\vec{b}|=\sqrt{41}=\sqrt{41}$.}]{Cho $\vec{a}=(-1;1;-2)$, $\vec{b}=(1;-5;-1)$. Tính $|-\vec{a}+\vec{b}|^2$.}{41}

\cauTLN[id=C2-B8-D07-TH-075,muc=TH,loigiai={Ta có $-2\vec{a}-\vec{b}=(9;-1;-7)$ nên $|-2\vec{a}-\vec{b}|=\sqrt{131}=\sqrt{131}$.}]{Cho $\vec{a}=(-4;1;5)$, $\vec{b}=(-1;-1;-3)$. Tính $|-2\vec{a}-\vec{b}|^2$.}{131}

\cauTLN[id=C2-B8-D07-TH-076,muc=TH,loigiai={Tính từng thành phần: $\vec{u}=-2\vec{a}+\vec{b}+2\vec{c}=(0;1;0)$.}]{Trong không gian $Oxyz$, cho $\vec{a}=(-1;-5;-4)$, $\vec{b}=(2;-1;0)$, $\vec{c}=(-2;-4;-4)$. Tính tổng các tọa độ của vectơ $\vec{u}=-2\vec{a}+\vec{b}+2\vec{c}$.}{1}

\cauTLN[id=C2-B8-D07-TH-077,muc=TH,loigiai={Tính từng thành phần: $\vec{u}=-\vec{a}-\vec{b}=(-4;4;5)$.}]{Trong không gian $Oxyz$, cho $\vec{a}=(3;1;-3)$, $\vec{b}=(1;-5;-2)$. Tính tổng các tọa độ của vectơ $\vec{u}=-\vec{a}-\vec{b}$.}{5}

\cauTLN[id=C2-B8-D07-TH-078,muc=TH,loigiai={Tính từng thành phần: $\vec{u}=3\vec{a}-\vec{b}=(15;-8;-8)$.}]{Trong không gian $Oxyz$, cho $\vec{a}=(5;-1;-2)$, $\vec{b}=(0;5;2)$. Tính tổng các tọa độ của vectơ $\vec{u}=3\vec{a}-\vec{b}$.}{-1}

\cauTLN[id=C2-B8-D07-TH-079,muc=TH,loigiai={Tính từng thành phần: $\vec{u}=2\vec{a}+2\vec{b}=(-14;4;-6)$.}]{Trong không gian $Oxyz$, cho $\vec{a}=(-2;2;2)$, $\vec{b}=(-5;0;-5)$. Tính tổng các tọa độ của vectơ $\vec{u}=2\vec{a}+2\vec{b}$.}{-16}

\cauTLN[id=C2-B8-D07-TH-080,muc=TH,loigiai={Tính từng thành phần: $\vec{u}=-\vec{a}+2\vec{b}+\vec{c}=(-12;-10;3)$.}]{Trong không gian $Oxyz$, cho $\vec{a}=(3;4;-4)$, $\vec{b}=(-4;-4;0)$, $\vec{c}=(-1;2;-1)$. Tính tổng các tọa độ của vectơ $\vec{u}=-\vec{a}+2\vec{b}+\vec{c}$.}{-19}

\cauTN[id=C2-B8-D07-VD-001,muc=VD,loigiai={Giả sử $\vec{c}=x\vec{a}+y\vec{b}$ ta có hệ $\begin{cases}12=3x+3y\\1=x+0y\\8=-x+3y\end{cases}$, giải ra $x=1$, $y=3$.}]{Cho $\vec{a}=(3;1;-1)$, $\vec{b}=(3;0;3)$, $\vec{c}=(12;1;8)$. Biết $\vec{c}=x\vec{a}+y\vec{b}$. Cặp $(x;y)$ là}
  {$(1;3)$.}
  {$(3;1)$.}{$(-1;3)$.}{$(1;-3)$.}

\cauTN[id=C2-B8-D07-VD-002,muc=VD,loigiai={Giả sử $\vec{c}=x\vec{a}+y\vec{b}$ ta có hệ $\begin{cases}2=2x+2y\\-12=3x-2y\\2=-x+0y\end{cases}$, giải ra $x=-2$, $y=3$.}]{Cho $\vec{a}=(2;3;-1)$, $\vec{b}=(2;-2;0)$, $\vec{c}=(2;-12;2)$. Biết $\vec{c}=x\vec{a}+y\vec{b}$. Cặp $(x;y)$ là}
  {$(-2;3)$.}
  {$(3;-2)$.}{$(2;3)$.}{$(-2;-3)$.}

\cauTN[id=C2-B8-D07-VD-003,muc=VD,loigiai={Giả sử $\vec{c}=x\vec{a}+y\vec{b}$ ta có hệ $\begin{cases}6=x+y\\-18=-3x-3y\\9=3x+0y\end{cases}$, giải ra $x=3$, $y=3$.}]{Cho $\vec{a}=(1;-3;3)$, $\vec{b}=(1;-3;0)$, $\vec{c}=(6;-18;9)$. Biết $\vec{c}=x\vec{a}+y\vec{b}$. Cặp $(x;y)$ là}
  {$(3;3)$.}
  {$(-3;3)$.}{$(3;-3)$.}{$(4;3)$.}

\cauTN[id=C2-B8-D07-VD-004,muc=VD,loigiai={Giả sử $\vec{c}=x\vec{a}+y\vec{b}$ ta có hệ $\begin{cases}12=x+3y\\-6=0x-2y\\6=-x+3y\end{cases}$, giải ra $x=3$, $y=3$.}]{Cho $\vec{a}=(1;0;-1)$, $\vec{b}=(3;-2;3)$, $\vec{c}=(12;-6;6)$. Biết $\vec{c}=x\vec{a}+y\vec{b}$. Cặp $(x;y)$ là}
  {$(3;3)$.}
  {$(-3;3)$.}{$(3;-3)$.}{$(4;3)$.}

\cauTN[id=C2-B8-D07-VD-005,muc=VD,loigiai={Giả sử $\vec{c}=x\vec{a}+y\vec{b}$ ta có hệ $\begin{cases}5=x-y\\-4=-2x-y\\1=x+y\end{cases}$, giải ra $x=3$, $y=-2$.}]{Cho $\vec{a}=(1;-2;1)$, $\vec{b}=(-1;-1;1)$, $\vec{c}=(5;-4;1)$. Biết $\vec{c}=x\vec{a}+y\vec{b}$. Cặp $(x;y)$ là}
  {$(3;-2)$.}
  {$(-2;3)$.}{$(-3;-2)$.}{$(3;2)$.}

\cauTN[id=C2-B8-D07-VD-006,muc=VD,loigiai={Giả sử $\vec{c}=x\vec{a}+y\vec{b}$ ta có hệ $\begin{cases}1=-x+y\\-5=-x-y\\4=2x+0y\end{cases}$, giải ra $x=2$, $y=3$.}]{Cho $\vec{a}=(-1;-1;2)$, $\vec{b}=(1;-1;0)$, $\vec{c}=(1;-5;4)$. Biết $\vec{c}=x\vec{a}+y\vec{b}$. Cặp $(x;y)$ là}
  {$(2;3)$.}
  {$(3;2)$.}{$(-2;3)$.}{$(2;-3)$.}

\cauTN[id=C2-B8-D07-VD-007,muc=VD,loigiai={Giả sử $\vec{c}=x\vec{a}+y\vec{b}$ ta có hệ $\begin{cases}-4=0x+2y\\6=0x-3y\\-6=-x+2y\end{cases}$, giải ra $x=2$, $y=-2$.}]{Cho $\vec{a}=(0;0;-1)$, $\vec{b}=(2;-3;2)$, $\vec{c}=(-4;6;-6)$. Biết $\vec{c}=x\vec{a}+y\vec{b}$. Cặp $(x;y)$ là}
  {$(2;-2)$.}
  {$(-2;2)$.}{$(-2;-2)$.}{$(2;2)$.}

\cauTN[id=C2-B8-D07-VD-008,muc=VD,loigiai={Giả sử $\vec{c}=x\vec{a}+y\vec{b}$ ta có hệ $\begin{cases}3=x+y\\4=x+2y\\0=-x+2y\end{cases}$, giải ra $x=2$, $y=1$.}]{Cho $\vec{a}=(1;1;-1)$, $\vec{b}=(1;2;2)$, $\vec{c}=(3;4;0)$. Biết $\vec{c}=x\vec{a}+y\vec{b}$. Cặp $(x;y)$ là}
  {$(2;1)$.}
  {$(1;2)$.}{$(-2;1)$.}{$(2;-1)$.}

\cauTN[id=C2-B8-D07-VD-009,muc=VD,loigiai={Giả sử $\vec{c}=x\vec{a}+y\vec{b}$ ta có hệ $\begin{cases}4=-x-y\\4=-x-y\\-2=2x-y\end{cases}$, giải ra $x=-2$, $y=-2$.}]{Cho $\vec{a}=(-1;-1;2)$, $\vec{b}=(-1;-1;-1)$, $\vec{c}=(4;4;-2)$. Biết $\vec{c}=x\vec{a}+y\vec{b}$. Cặp $(x;y)$ là}
  {$(-2;-2)$.}
  {$(2;-2)$.}{$(-2;2)$.}{$(-1;-2)$.}

\cauTN[id=C2-B8-D07-VD-010,muc=VD,loigiai={Giả sử $\vec{c}=x\vec{a}+y\vec{b}$ ta có hệ $\begin{cases}8=-3x-2y\\1=-2x+3y\\1=-2x+3y\end{cases}$, giải ra $x=-2$, $y=-1$.}]{Cho $\vec{a}=(-3;-2;-2)$, $\vec{b}=(-2;3;3)$, $\vec{c}=(8;1;1)$. Biết $\vec{c}=x\vec{a}+y\vec{b}$. Cặp $(x;y)$ là}
  {$(-2;-1)$.}
  {$(-1;-2)$.}{$(2;-1)$.}{$(-2;1)$.}

\cauTN[id=C2-B8-D07-VD-011,muc=VD,loigiai={Giả sử $\vec{c}=x\vec{a}+y\vec{b}$ ta có hệ $\begin{cases}4=-2x+y\\4=2x+3y\\-4=0x-2y\end{cases}$, giải ra $x=-1$, $y=2$.}]{Cho $\vec{a}=(-2;2;0)$, $\vec{b}=(1;3;-2)$, $\vec{c}=(4;4;-4)$. Biết $\vec{c}=x\vec{a}+y\vec{b}$. Cặp $(x;y)$ là}
  {$(-1;2)$.}
  {$(2;-1)$.}{$(1;2)$.}{$(-1;-2)$.}

\cauTN[id=C2-B8-D07-VD-012,muc=VD,loigiai={Giả sử $\vec{c}=x\vec{a}+y\vec{b}$ ta có hệ $\begin{cases}3=-x+2y\\3=-3x+0y\\3=0x+3y\end{cases}$, giải ra $x=-1$, $y=1$.}]{Cho $\vec{a}=(-1;-3;0)$, $\vec{b}=(2;0;3)$, $\vec{c}=(3;3;3)$. Biết $\vec{c}=x\vec{a}+y\vec{b}$. Cặp $(x;y)$ là}
  {$(-1;1)$.}
  {$(1;-1)$.}{$(1;1)$.}{$(-1;-1)$.}

\cauTN[id=C2-B8-D07-VD-013,muc=VD,loigiai={Ta có $-2\vec{a}+2\vec{b}=(-2;-2;18)$ nên $|-2\vec{a}+2\vec{b}|=\sqrt{332}=2\sqrt{83}$.}]{Cho $\vec{a}=(-2;1;-5)$, $\vec{b}=(-3;0;4)$. Độ dài của vectơ $-2\vec{a}+2\vec{b}$ bằng}
  {$2\sqrt{83}$.}
  {$3\sqrt{37}$.}{$\sqrt{331}$.}{$4\sqrt{21}$.}

\cauTN[id=C2-B8-D07-VD-014,muc=VD,loigiai={Ta có $3\vec{a}+\vec{b}=(-16;-16;-13)$ nên $|3\vec{a}+\vec{b}|=\sqrt{681}=\sqrt{681}$.}]{Cho $\vec{a}=(-6;-4;-4)$, $\vec{b}=(2;-4;-1)$. Độ dài của vectơ $3\vec{a}+\vec{b}$ bằng}
  {$\sqrt{681}$.}
  {$\sqrt{682}$.}{$2\sqrt{170}$.}{$\sqrt{685}$.}

\cauTN[id=C2-B8-D07-VD-015,muc=VD,loigiai={Ta có $3\vec{a}+2\vec{b}=(15;-18;-8)$ nên $|3\vec{a}+2\vec{b}|=\sqrt{613}=\sqrt{613}$.}]{Cho $\vec{a}=(5;-4;0)$, $\vec{b}=(0;-3;-4)$. Độ dài của vectơ $3\vec{a}+2\vec{b}$ bằng}
  {$\sqrt{613}$.}
  {$\sqrt{614}$.}{$6\sqrt{17}$.}{$\sqrt{617}$.}

\cauTN[id=C2-B8-D07-VD-016,muc=VD,loigiai={Ta có $2\vec{a}-3\vec{b}=(-3;-21;9)$ nên $|2\vec{a}-3\vec{b}|=\sqrt{531}=3\sqrt{59}$.}]{Cho $\vec{a}=(3;-3;6)$, $\vec{b}=(3;5;1)$. Độ dài của vectơ $2\vec{a}-3\vec{b}$ bằng}
  {$3\sqrt{59}$.}
  {$2\sqrt{133}$.}{$\sqrt{530}$.}{$\sqrt{535}$.}

\cauTN[id=C2-B8-D07-VD-017,muc=VD,loigiai={Ta có $3\vec{a}+\vec{b}=(-1;-18;-3)$ nên $|3\vec{a}+\vec{b}|=\sqrt{334}=\sqrt{334}$.}]{Cho $\vec{a}=(-1;-5;-3)$, $\vec{b}=(2;-3;6)$. Độ dài của vectơ $3\vec{a}+\vec{b}$ bằng}
  {$\sqrt{334}$.}
  {$\sqrt{335}$.}{$3\sqrt{37}$.}{$13\sqrt{2}$.}

\cauTN[id=C2-B8-D07-VD-018,muc=VD,loigiai={Ta có $-\vec{a}-\vec{b}=(-6;6;-1)$ nên $|-\vec{a}-\vec{b}|=\sqrt{73}=\sqrt{73}$.}]{Cho $\vec{a}=(4;-5;-3)$, $\vec{b}=(2;-1;4)$. Độ dài của vectơ $-\vec{a}-\vec{b}$ bằng}
  {$\sqrt{73}$.}
  {$\sqrt{74}$.}{$6\sqrt{2}$.}{$\sqrt{77}$.}

\cauTN[id=C2-B8-D07-VD-019,muc=VD,loigiai={Ta có $2\vec{a}+3\vec{b}=(20;-14;14)$ nên $|2\vec{a}+3\vec{b}|=\sqrt{792}=6\sqrt{22}$.}]{Cho $\vec{a}=(1;-4;-2)$, $\vec{b}=(6;-2;6)$. Độ dài của vectơ $2\vec{a}+3\vec{b}$ bằng}
  {$6\sqrt{22}$.}
  {$\sqrt{793}$.}{$\sqrt{791}$.}{$2\sqrt{199}$.}

\cauTN[id=C2-B8-D07-VD-020,muc=VD,loigiai={Ta có $2\vec{a}-\vec{b}=(5;-4;7)$ nên $|2\vec{a}-\vec{b}|=\sqrt{90}=3\sqrt{10}$.}]{Cho $\vec{a}=(5;-4;6)$, $\vec{b}=(5;-4;5)$. Độ dài của vectơ $2\vec{a}-\vec{b}$ bằng}
  {$3\sqrt{10}$.}
  {$\sqrt{91}$.}{$\sqrt{89}$.}{$\sqrt{94}$.}

\cauTN[id=C2-B8-D07-VD-021,muc=VD,loigiai={Ta có $-2\vec{a}-3\vec{b}=(-3;7;-2)$ nên $|-2\vec{a}-3\vec{b}|=\sqrt{62}=\sqrt{62}$.}]{Cho $\vec{a}=(-3;4;1)$, $\vec{b}=(3;-5;0)$. Độ dài của vectơ $-2\vec{a}-3\vec{b}$ bằng}
  {$\sqrt{62}$.}
  {$3\sqrt{7}$.}{$\sqrt{61}$.}{$\sqrt{66}$.}

\cauTN[id=C2-B8-D07-VD-022,muc=VD,loigiai={Ta có $3\vec{a}+3\vec{b}=(0;-18;0)$ nên $|3\vec{a}+3\vec{b}|=\sqrt{324}=18$.}]{Cho $\vec{a}=(-5;-2;0)$, $\vec{b}=(5;-4;0)$. Độ dài của vectơ $3\vec{a}+3\vec{b}$ bằng}
  {$18$.}
  {$5\sqrt{13}$.}{$\sqrt{323}$.}{$2\sqrt{82}$.}

\cauTN[id=C2-B8-D07-VD-023,muc=VD,loigiai={Ta có $2\vec{a}+3\vec{b}=(10;-10;16)$ nên $|2\vec{a}+3\vec{b}|=\sqrt{456}=2\sqrt{114}$.}]{Cho $\vec{a}=(5;-5;5)$, $\vec{b}=(0;0;2)$. Độ dài của vectơ $2\vec{a}+3\vec{b}$ bằng}
  {$2\sqrt{114}$.}
  {$\sqrt{457}$.}{$\sqrt{455}$.}{$2\sqrt{115}$.}

\cauTN[id=C2-B8-D07-VD-024,muc=VD,loigiai={Ta có $3\vec{a}-\vec{b}=(-10;-6;3)$ nên $|3\vec{a}-\vec{b}|=\sqrt{145}=\sqrt{145}$.}]{Cho $\vec{a}=(-3;-3;2)$, $\vec{b}=(1;-3;3)$. Độ dài của vectơ $3\vec{a}-\vec{b}$ bằng}
  {$\sqrt{145}$.}
  {$\sqrt{146}$.}{$12$.}{$\sqrt{149}$.}

\cauTN[id=C2-B8-D07-VD-025,muc=VD,loigiai={Ta có $B=A+\overrightarrow{AB}=(3;3;4)+((11;-6;-5))=(14;-3;-1)$.}]{Cho điểm $A(3;3;4)$, $\vec{a}=(-5;1;6)$, $\vec{b}=(-6;5;-1)$. Điểm $B$ thoả mãn $\overrightarrow{AB}=-1\vec{a}-\vec{b}$ có tọa độ là}
  {$(14;-3;-1)$.}
  {$(14;-3;0)$.}{$(14;-3;1)$.}{$(16;-3;-1)$.}

\cauTN[id=C2-B8-D07-VD-026,muc=VD,loigiai={Ta có $B=A+\overrightarrow{AB}=(4;6;-1)+((-1;-4;8))=(3;2;7)$.}]{Cho điểm $A(4;6;-1)$, $\vec{a}=(-4;-5;2)$, $\vec{b}=(3;1;6)$. Điểm $B$ thoả mãn $\overrightarrow{AB}=\vec{a}+\vec{b}$ có tọa độ là}
  {$(3;2;7)$.}
  {$(3;2;9)$.}{$(-3;-2;-7)$.}{$(3;2;8)$.}

\cauTN[id=C2-B8-D07-VD-027,muc=VD,loigiai={Ta có $B=A+\overrightarrow{AB}=(6;5;5)+((-15;-3;3))=(-9;2;8)$.}]{Cho điểm $A(6;5;5)$, $\vec{a}=(-6;1;2)$, $\vec{b}=(3;-6;-3)$. Điểm $B$ thoả mãn $\overrightarrow{AB}=3\vec{a}+\vec{b}$ có tọa độ là}
  {$(-9;2;8)$.}
  {$(-9;2;-8)$.}{$(-9;-2;8)$.}{$(-7;2;8)$.}

\cauTN[id=C2-B8-D07-VD-028,muc=VD,loigiai={Ta có $B=A+\overrightarrow{AB}=(0;5;-6)+((12;7;7))=(12;12;1)$.}]{Cho điểm $A(0;5;-6)$, $\vec{a}=(6;6;3)$, $\vec{b}=(0;5;-1)$. Điểm $B$ thoả mãn $\overrightarrow{AB}=2\vec{a}-\vec{b}$ có tọa độ là}
  {$(12;12;1)$.}
  {$(-12;12;1)$.}{$(12;11;1)$.}{$(1;12;12)$.}

\cauTN[id=C2-B8-D07-VD-029,muc=VD,loigiai={Ta có $B=A+\overrightarrow{AB}=(-3;-3;-3)+((-6;4;0))=(-9;1;-3)$.}]{Cho điểm $A(-3;-3;-3)$, $\vec{a}=(-4;6;-2)$, $\vec{b}=(5;-5;1)$. Điểm $B$ thoả mãn $\overrightarrow{AB}=-1\vec{a}-2\vec{b}$ có tọa độ là}
  {$(-9;1;-3)$.}
  {$(-10;1;-3)$.}{$(-9;1;3)$.}{$(9;-1;3)$.}

\cauTN[id=C2-B8-D07-VD-030,muc=VD,loigiai={Ta có $B=A+\overrightarrow{AB}=(-6;4;-5)+((7;4;-5))=(1;8;-10)$.}]{Cho điểm $A(-6;4;-5)$, $\vec{a}=(-6;-6;-1)$, $\vec{b}=(-1;2;6)$. Điểm $B$ thoả mãn $\overrightarrow{AB}=-1\vec{a}-\vec{b}$ có tọa độ là}
  {$(1;8;-10)$.}
  {$(-13;0;0)$.}{$(-10;8;1)$.}{$(1;8;-11)$.}

\cauTN[id=C2-B8-D07-VD-031,muc=VD,loigiai={Ta có $B=A+\overrightarrow{AB}=(-6;5;0)+((20;18;-22))=(14;23;-22)$.}]{Cho điểm $A(-6;5;0)$, $\vec{a}=(4;4;-6)$, $\vec{b}=(-4;-3;2)$. Điểm $B$ thoả mãn $\overrightarrow{AB}=3\vec{a}-2\vec{b}$ có tọa độ là}
  {$(14;23;-22)$.}
  {$(-22;23;14)$.}{$(14;-22;23)$.}{$(14;23;22)$.}

\cauTN[id=C2-B8-D07-VD-032,muc=VD,loigiai={Ta có $B=A+\overrightarrow{AB}=(6;-6;-3)+((-4;-8;12))=(2;-14;9)$.}]{Cho điểm $A(6;-6;-3)$, $\vec{a}=(-2;-1;1)$, $\vec{b}=(0;3;-5)$. Điểm $B$ thoả mãn $\overrightarrow{AB}=2\vec{a}-2\vec{b}$ có tọa độ là}
  {$(2;-14;9)$.}
  {$(3;-14;9)$.}{$(2;-16;9)$.}{$(-2;-14;9)$.}

\cauTN[id=C2-B8-D07-VD-033,muc=VD,loigiai={Ta có $B=A+\overrightarrow{AB}=(-2;1;-3)+((1;5;-2))=(-1;6;-5)$.}]{Cho điểm $A(-2;1;-3)$, $\vec{a}=(3;-1;1)$, $\vec{b}=(-4;-4;1)$. Điểm $B$ thoả mãn $\overrightarrow{AB}=-1\vec{a}-\vec{b}$ có tọa độ là}
  {$(-1;6;-5)$.}
  {$(-1;6;5)$.}{$(1;6;-5)$.}{$(-1;6;-7)$.}

\cauTN[id=C2-B8-D07-VD-034,muc=VD,loigiai={Ta có $B=A+\overrightarrow{AB}=(-5;-4;-1)+((-2;3;-10))=(-7;-1;-11)$.}]{Cho điểm $A(-5;-4;-1)$, $\vec{a}=(0;6;-6)$, $\vec{b}=(2;3;4)$. Điểm $B$ thoả mãn $\overrightarrow{AB}=\vec{a}-\vec{b}$ có tọa độ là}
  {$(-7;-1;-11)$.}
  {$(-7;1;-11)$.}{$(-1;-7;-11)$.}{$(-7;-1;11)$.}

\cauTN[id=C2-B8-D07-VD-035,muc=VD,loigiai={Ta có $B=A+\overrightarrow{AB}=(4;-4;-6)+((12;-1;-2))=(16;-5;-8)$.}]{Cho điểm $A(4;-4;-6)$, $\vec{a}=(4;-1;0)$, $\vec{b}=(4;1;-2)$. Điểm $B$ thoả mãn $\overrightarrow{AB}=2\vec{a}+\vec{b}$ có tọa độ là}
  {$(16;-5;-8)$.}
  {$(16;-6;-8)$.}{$(16;-3;-8)$.}{$(16;-5;-10)$.}

\cauTN[id=C2-B8-D07-VD-036,muc=VD,loigiai={Ta có $B=A+\overrightarrow{AB}=(-1;-4;0)+((-8;0;-2))=(-9;-4;-2)$.}]{Cho điểm $A(-1;-4;0)$, $\vec{a}=(-4;6;-6)$, $\vec{b}=(-2;-3;2)$. Điểm $B$ thoả mãn $\overrightarrow{AB}=\vec{a}+2\vec{b}$ có tọa độ là}
  {$(-9;-4;-2)$.}
  {$(-9;-4;2)$.}{$(-9;-4;-1)$.}{$(9;4;2)$.}

\cauDS[id=C2-B8-D07-VD-037,muc=VD,loigiai={Kiểm tra từng mệnh đề bằng công thức tọa độ. a) Đúng. b) Đúng. c) Sai. d) Đúng.}]{Trong không gian $Oxyz$, cho $A(-2;3;-4)$, $B(-2;3;1)$.  Cho $\vec{a}=(3;3;6)$, $\vec{b}=(2;4;-1)$, $\vec{c}=(-5;1;1)$. Xét tính đúng sai của các mệnh đề sau:}{%
  \y{Độ dài của vectơ $\vec{a}$ bằng $3\sqrt{6}$.}{d}
  \y{Độ dài đoạn thẳng $AB$ bằng $5$.}{d}
  \y{Vectơ $\vec{u}=-\vec{a}+\vec{b}$ có tọa độ $(-1;2;-7)$.}{s}
  \y{Với $\vec{a}=(3;3;6)$, $\vec{b}=(2;4;-1)$ thì $|2\vec{a}+3\vec{b}|=3\sqrt{61}$.}{d}}

\cauDS[id=C2-B8-D07-VD-038,muc=VD,loigiai={Kiểm tra từng mệnh đề bằng công thức tọa độ. a) Đúng. b) Đúng. c) Sai. d) Sai.}]{Trong không gian $Oxyz$, cho $A(-5;-2;-2)$, $B(-1;-2;0)$.  Cho $\vec{a}=(-1;6;-4)$. Xét tính đúng sai của các mệnh đề sau:}{%
  \y{Độ dài đoạn thẳng $AB$ bằng $2\sqrt{5}$.}{d}
  \y{Độ dài của vectơ $\vec{a}$ bằng $\sqrt{53}$.}{d}
  \y{Với $\vec{a}=(-1;6;-4)$, $\vec{b}=(-5;-2;-2)$ thì $|3\vec{a}+3\vec{b}|=\sqrt{793}$.}{s}
  \y{Độ dài của vectơ $\vec{a}$ bằng $2\sqrt{13}$.}{s}}

\cauDS[id=C2-B8-D07-VD-039,muc=VD,loigiai={Kiểm tra từng mệnh đề bằng công thức tọa độ. a) Sai. b) Đúng. c) Đúng. d) Sai.}]{Trong không gian $Oxyz$, cho $A(0;2;-1)$, $B(-3;-1;-4)$.  Cho $\vec{a}=(6;-5;-6)$. Xét tính đúng sai của các mệnh đề sau:}{%
  \y{Với $\vec{a}=(6;-5;-6)$, $\vec{b}=(-4;4;-6)$ thì $|3\vec{a}+\vec{b}|=2\sqrt{223}$.}{s}
  \y{Độ dài của vectơ $\vec{a}$ bằng $\sqrt{97}$.}{d}
  \y{Với $\vec{a}=(6;-5;-6)$, $\vec{b}=(-4;4;-6)$ thì $|2\vec{a}-3\vec{b}|=2\sqrt{274}$.}{d}
  \y{Độ dài đoạn thẳng $AB$ bằng $\sqrt{31}$.}{s}}

\cauDS[id=C2-B8-D07-VD-040,muc=VD,loigiai={Kiểm tra từng mệnh đề bằng công thức tọa độ. a) Đúng. b) Sai. c) Sai. d) Đúng.}]{Trong không gian $Oxyz$, cho $A(-5;3;-2)$, $B(3;6;1)$.  Cho $\vec{a}=(-1;-5;-1)$, $\vec{b}=(5;0;0)$, $\vec{c}=(1;-5;-3)$. Xét tính đúng sai của các mệnh đề sau:}{%
  \y{Độ dài đoạn thẳng $AB$ bằng $\sqrt{82}$.}{d}
  \y{Với $\vec{a}=(-1;-5;-1)$, $\vec{b}=(5;0;0)$ thì $|-2\vec{a}+\vec{b}|=2\sqrt{38}$.}{s}
  \y{Độ dài của vectơ $\vec{a}$ bằng $\sqrt{31}$.}{s}
  \y{Vectơ $\vec{u}=-\vec{a}+3\vec{b}$ có tọa độ $(16;5;1)$.}{d}}

\cauDS[id=C2-B8-D07-VD-041,muc=VD,loigiai={Kiểm tra từng mệnh đề bằng công thức tọa độ. a) Đúng. b) Đúng. c) Sai. d) Sai.}]{Trong không gian $Oxyz$, cho $A(3;5;5)$, $B(-3;-2;3)$.  Cho $\vec{a}=(6;-1;-1)$, $\vec{b}=(1;5;5)$, $\vec{c}=(2;-1;-5)$. Xét tính đúng sai của các mệnh đề sau:}{%
  \y{Vectơ $\vec{u}=3\vec{a}-3\vec{b}$ có tọa độ $(15;-18;-18)$.}{d}
  \y{Độ dài của vectơ $\vec{a}$ bằng $\sqrt{38}$.}{d}
  \y{Độ dài của vectơ $\vec{a}$ bằng $\sqrt{37}$.}{s}
  \y{Độ dài đoạn thẳng $AB$ bằng $3\sqrt{10}$.}{s}}

\cauDS[id=C2-B8-D07-VD-042,muc=VD,loigiai={Kiểm tra từng mệnh đề bằng công thức tọa độ. a) Đúng. b) Đúng. c) Sai. d) Sai.}]{Trong không gian $Oxyz$, cho $A(0;2;-2)$, $B(3;-2;-5)$.  Cho $\vec{a}=(1;0;6)$, $\vec{b}=(-1;0;-3)$, $\vec{c}=(-6;2;3)$. Xét tính đúng sai của các mệnh đề sau:}{%
  \y{Với $\vec{a}=(1;0;6)$, $\vec{b}=(-1;0;-3)$ thì $|2\vec{a}-\vec{b}|=3\sqrt{26}$.}{d}
  \y{Với $\vec{a}=(1;0;6)$, $\vec{b}=(-1;0;-3)$ thì $|-2\vec{a}+\vec{b}|=3\sqrt{26}$.}{d}
  \y{Độ dài đoạn thẳng $AB$ bằng $\sqrt{35}$.}{s}
  \y{Vectơ $\vec{u}=-\vec{a}+2\vec{b}$ có tọa độ $(-3;0;-11)$.}{s}}

\cauDS[id=C2-B8-D07-VD-043,muc=VD,loigiai={Kiểm tra từng mệnh đề bằng công thức tọa độ. a) Sai. b) Sai. c) Sai. d) Đúng.}]{Trong không gian $Oxyz$, cho $A(4;-6;-1)$, $B(4;-2;3)$.  Cho $\vec{a}=(-3;2;-1)$, $\vec{b}=(-5;-6;-5)$, $\vec{c}=(-2;6;2)$. Xét tính đúng sai của các mệnh đề sau:}{%
  \y{Độ dài của vectơ $\vec{a}$ bằng $3\sqrt{2}$.}{s}
  \y{Độ dài của vectơ $\vec{a}$ bằng $\sqrt{13}$.}{s}
  \y{Vectơ $\vec{u}=2\vec{a}+\vec{b}$ có tọa độ $(-11;-2;7)$.}{s}
  \y{Độ dài đoạn thẳng $AB$ bằng $4\sqrt{2}$.}{d}}

\cauDS[id=C2-B8-D07-VD-044,muc=VD,loigiai={Kiểm tra từng mệnh đề bằng công thức tọa độ. a) Đúng. b) Đúng. c) Đúng. d) Sai.}]{Trong không gian $Oxyz$, cho $A(-6;-5;-1)$, $B(1;4;-1)$.  Cho $\vec{a}=(6;-2;-6)$, $\vec{b}=(0;-3;6)$, $\vec{c}=(5;2;0)$. Xét tính đúng sai của các mệnh đề sau:}{%
  \y{Độ dài đoạn thẳng $AB$ bằng $\sqrt{130}$.}{d}
  \y{Vectơ $\vec{u}=-\vec{a}-3\vec{b}$ có tọa độ $(-6;11;-12)$.}{d}
  \y{Độ dài của vectơ $\vec{a}$ bằng $2\sqrt{19}$.}{d}
  \y{Độ dài của vectơ $\vec{a}$ bằng $4\sqrt{5}$.}{s}}

\cauDS[id=C2-B8-D07-VD-045,muc=VD,loigiai={Kiểm tra từng mệnh đề bằng công thức tọa độ. a) Đúng. b) Đúng. c) Sai. d) Sai.}]{Trong không gian $Oxyz$, cho $A(-1;-2;1)$, $B(-2;2;-4)$.  Cho $\vec{a}=(-4;-5;-4)$. Xét tính đúng sai của các mệnh đề sau:}{%
  \y{Với $\vec{a}=(-4;-5;-4)$, $\vec{b}=(-2;-4;5)$ thì $|2\vec{a}-3\vec{b}|=\sqrt{537}$.}{d}
  \y{Với $\vec{a}=(-4;-5;-4)$, $\vec{b}=(-2;-4;5)$ thì $|-2\vec{a}+\vec{b}|=\sqrt{241}$.}{d}
  \y{Độ dài của vectơ $\vec{a}$ bằng $2\sqrt{14}$.}{s}
  \y{Độ dài đoạn thẳng $AB$ bằng $\sqrt{43}$.}{s}}

\cauDS[id=C2-B8-D07-VD-046,muc=VD,loigiai={Kiểm tra từng mệnh đề bằng công thức tọa độ. a) Sai. b) Đúng. c) Đúng. d) Sai.}]{Trong không gian $Oxyz$, cho $\vec{a}=(5;0;4)$, $\vec{b}=(-1;-6;3)$, $\vec{c}=(6;2;1)$. Xét tính đúng sai của các mệnh đề sau:}{%
  \y{Độ dài của vectơ $\vec{a}$ bằng $\sqrt{42}$.}{s}
  \y{Vectơ $\vec{u}=-2\vec{a}+3\vec{b}+2\vec{c}$ có tọa độ $(-1;-14;3)$.}{d}
  \y{Với $\vec{a}=(5;0;4)$, $\vec{b}=(-1;-6;3)$ thì $|3\vec{a}-3\vec{b}|=3\sqrt{73}$.}{d}
  \y{Vectơ $\vec{u}=3\vec{a}+3\vec{b}$ có tọa độ $(12;-18;-21)$.}{s}}

\cauDS[id=C2-B8-D07-VD-047,muc=VD,loigiai={Kiểm tra từng mệnh đề bằng công thức tọa độ. a) Sai. b) Sai. c) Sai. d) Đúng.}]{Trong không gian $Oxyz$, cho $A(-2;-2;-3)$, $B(-2;-1;6)$.  Cho $\vec{a}=(4;-3;-2)$, $\vec{b}=(4;-6;-4)$, $\vec{c}=(0;-2;3)$. Xét tính đúng sai của các mệnh đề sau:}{%
  \y{Độ dài đoạn thẳng $AB$ bằng $\sqrt{86}$.}{s}
  \y{Vectơ $\vec{u}=-2\vec{a}+3\vec{b}-\vec{c}$ có tọa độ $(4;-11;-11)$.}{s}
  \y{Với $\vec{a}=(4;-3;-2)$, $\vec{b}=(4;-6;-4)$ thì $|-\vec{a}-3\vec{b}|=2\sqrt{223}$.}{s}
  \y{Với $\vec{a}=(4;-3;-2)$, $\vec{b}=(4;-6;-4)$ thì $|3\vec{a}+\vec{b}|=\sqrt{581}$.}{d}}

\cauDS[id=C2-B8-D07-VD-048,muc=VD,loigiai={Kiểm tra từng mệnh đề bằng công thức tọa độ. a) Đúng. b) Sai. c) Sai. d) Đúng.}]{Trong không gian $Oxyz$, cho $A(-1;3;-3)$, $B(-3;-4;-6)$.  Cho $\vec{a}=(0;6;4)$. Xét tính đúng sai của các mệnh đề sau:}{%
  \y{Độ dài của vectơ $\vec{a}$ bằng $2\sqrt{13}$.}{d}
  \y{Độ dài của vectơ $\vec{a}$ bằng $\sqrt{53}$.}{s}
  \y{Độ dài đoạn thẳng $AB$ bằng $\sqrt{61}$.}{s}
  \y{Với $\vec{a}=(0;6;4)$, $\vec{b}=(-6;-4;6)$ thì $|-2\vec{a}+3\vec{b}|=10\sqrt{10}$.}{d}}

\cauTLN[id=C2-B8-D07-VD-049,muc=VD,loigiai={Giả sử $\vec{c}=x\vec{a}+y\vec{b}$ ta có hệ $\begin{cases}6=3x+3y\\5=-2x+y\\6=0x+2y\end{cases}$, giải ra $x=-1$, $y=3$.}]{Cho $\vec{a}=(3;-2;0)$, $\vec{b}=(3;1;2)$, $\vec{c}=(6;5;6)$. Biết $\vec{c}=x\vec{a}+y\vec{b}$. Tính $x+2y$.}{5}

\cauTLN[id=C2-B8-D07-VD-050,muc=VD,loigiai={Giả sử $\vec{c}=x\vec{a}+y\vec{b}$ ta có hệ $\begin{cases}-11=x-3y\\3=-3x-y\\-15=3x-3y\end{cases}$, giải ra $x=-2$, $y=3$.}]{Cho $\vec{a}=(1;-3;3)$, $\vec{b}=(-3;-1;-3)$, $\vec{c}=(-11;3;-15)$. Biết $\vec{c}=x\vec{a}+y\vec{b}$. Tính $x+2y$.}{4}

\cauTLN[id=C2-B8-D07-VD-051,muc=VD,loigiai={Giả sử $\vec{c}=x\vec{a}+y\vec{b}$ ta có hệ $\begin{cases}0=3x-3y\\-6=-x-2y\\10=2x+3y\end{cases}$, giải ra $x=2$, $y=2$.}]{Cho $\vec{a}=(3;-1;2)$, $\vec{b}=(-3;-2;3)$, $\vec{c}=(0;-6;10)$. Biết $\vec{c}=x\vec{a}+y\vec{b}$. Tính $x+2y$.}{6}

\cauTLN[id=C2-B8-D07-VD-052,muc=VD,loigiai={Giả sử $\vec{c}=x\vec{a}+y\vec{b}$ ta có hệ $\begin{cases}4=-2x+0y\\1=-2x+3y\\8=-3x-2y\end{cases}$, giải ra $x=-2$, $y=-1$.}]{Cho $\vec{a}=(-2;-2;-3)$, $\vec{b}=(0;3;-2)$, $\vec{c}=(4;1;8)$. Biết $\vec{c}=x\vec{a}+y\vec{b}$. Tính $x+2y$.}{-4}

\cauTLN[id=C2-B8-D07-VD-053,muc=VD,loigiai={Ta có $2\vec{a}+3\vec{b}=(8;-7;6)$ nên $|2\vec{a}+3\vec{b}|=\sqrt{149}=\sqrt{149}$.}]{Cho $\vec{a}=(-2;1;6)$, $\vec{b}=(4;-3;-2)$. Tính $|2\vec{a}+3\vec{b}|^2$.}{149}

\cauTLN[id=C2-B8-D07-VD-054,muc=VD,loigiai={Ta có $2\vec{a}+3\vec{b}=(-16;8;-18)$ nên $|2\vec{a}+3\vec{b}|=\sqrt{644}=2\sqrt{161}$.}]{Cho $\vec{a}=(-2;4;0)$, $\vec{b}=(-4;0;-6)$. Tính $|2\vec{a}+3\vec{b}|^2$.}{644}

\cauTLN[id=C2-B8-D07-VD-055,muc=VD,loigiai={Ta có $3\vec{a}-\vec{b}=(9;21;9)$ nên $|3\vec{a}-\vec{b}|=\sqrt{603}=3\sqrt{67}$.}]{Cho $\vec{a}=(4;5;5)$, $\vec{b}=(3;-6;6)$. Tính $|3\vec{a}-\vec{b}|^2$.}{603}

\cauTLN[id=C2-B8-D07-VD-056,muc=VD,loigiai={Ta có $2\vec{a}+\vec{b}=(10;-7;-4)$ nên $|2\vec{a}+\vec{b}|=\sqrt{165}=\sqrt{165}$.}]{Cho $\vec{a}=(2;-6;0)$, $\vec{b}=(6;5;-4)$. Tính $|2\vec{a}+\vec{b}|^2$.}{165}

\cauTLN[id=C2-B8-D07-VD-057,muc=VD,loigiai={Ta có $B=A+\overrightarrow{AB}=(-5;4;2)+((-1;5;-4))=(-6;9;-2)$.}]{Cho điểm $A(-5;4;2)$, $\vec{a}=(1;-1;2)$, $\vec{b}=(0;-4;2)$. Điểm $B$ thoả mãn $\overrightarrow{AB}=-1\vec{a}-\vec{b}$. Tính $x_B+y_B+z_B$.}{1}

\cauTLN[id=C2-B8-D07-VD-058,muc=VD,loigiai={Ta có $B=A+\overrightarrow{AB}=(0;-1;2)+((8;2;-13))=(8;1;-11)$.}]{Cho điểm $A(0;-1;2)$, $\vec{a}=(-4;-4;-1)$, $\vec{b}=(6;3;-6)$. Điểm $B$ thoả mãn $\overrightarrow{AB}=\vec{a}+2\vec{b}$. Tính $x_B+y_B+z_B$.}{-2}

\cauTLN[id=C2-B8-D07-VD-059,muc=VD,loigiai={Ta có $B=A+\overrightarrow{AB}=(0;4;6)+((-14;-10;-8))=(-14;-6;-2)$.}]{Cho điểm $A(0;4;6)$, $\vec{a}=(-2;-4;2)$, $\vec{b}=(5;1;6)$. Điểm $B$ thoả mãn $\overrightarrow{AB}=2\vec{a}-2\vec{b}$. Tính $x_B+y_B+z_B$.}{-22}

\cauTLN[id=C2-B8-D07-VD-060,muc=VD,loigiai={Ta có $B=A+\overrightarrow{AB}=(5;3;-5)+((-1;19;-6))=(4;22;-11)$.}]{Cho điểm $A(5;3;-5)$, $\vec{a}=(-2;5;0)$, $\vec{b}=(-5;-4;6)$. Điểm $B$ thoả mãn $\overrightarrow{AB}=3\vec{a}-\vec{b}$. Tính $x_B+y_B+z_B$.}{15}
